\documentclass[a4paper,12pt]{refrep}
\usepackage{xcolor}

\usepackage{booktabs}
\usepackage{longtable}
\usepackage{rotating}
\usepackage{multicol}
\usepackage{multirow}
\usepackage[disable]{todonotes}
\usepackage[colorlinks=true]{hyperref}
\usepackage{gensymb}
\usepackage{makeidx}\makeindex
\makeatletter
\usepackage{fancyhdr}
\pagestyle{fancy}
\maxipagerulefalse
%
% Include XCSoar header and footer settings and used buttons 
\usepackage{calc}
\usepackage{fancyhdr}

\newcommand{\xcsoarheader}[1]{
\pagestyle{fancy}

% Add XCSoar User Manual title to the header
\fancyhead[L]{
\begin{tabular}{l}
\hspace*{-6cm} \em #1 \vspace*{2pt}
\end{tabular}
}

% Add page number to the footer (centered)
\fancyfoot{}
\fancyfoot[R]{\thepage}
}

% No line between content and header
\renewcommand{\headrulewidth}{0pt}

\fancypagestyle{plain}{
  % Clear fancy header and footer
  \fancyhf{}

  % Add page number to the footer (centered)
  \fancyfoot[R]{\thepage}

  % No line between content and header
  \renewcommand{\headrulewidth}{0pt}
}

\definecolor{buttongray}{rgb}{0.831,0.816,0.784}
\newcommand{\blink}[0]{$\triangleright$}
\newcommand{\bmenu}[1]{
	\fcolorbox {black}{buttongray}{{\sf{#1}}}
}
\newcommand{\bmenut}[2]{
	\fcolorbox {black}{buttongray}{
    \makebox[1.4cm][c]{
    	\begin{tabular}{c}
    	{\footnotesize\sf{#1}}\\
    	{\footnotesize\sf{#2}}
    	\end{tabular}
    }
  }
}
\newcommand{\bmenuth}[3]{
	\fcolorbox {black}{buttongray}{
    \makebox[1.4cm][c]{
    	\begin{tabular}{c}
    	{\footnotesize\sf{#1}}\\
    	{\footnotesize\sf{#2}}\\
    	{\footnotesize\sf{#3}}
    	\end{tabular}
    }
  }
}
\newcommand{\bmenus}[1]{
	\fcolorbox {black}{buttongray}{
    \makebox[1.4cm][c]{
      \begin{tabular}{c}
        {\footnotesize\sf{#1}}\\
    	  \\
      \end{tabular}
	  }
	}
}
\newcommand{\button}[1]{
	\fcolorbox {black}{buttongray}{{\sf #1}}
}

\newcommand{\infobox}[1]{
	\fcolorbox {black}{white}{\makebox[1.7cm][c]{\sf #1}}
}


\newenvironment{jspecs}{
\itemsep=2pt\topsep=3pt\partopsep=3pt\parskip=0pt
\begin{description}
\itemsep=2pt\topsep=3pt\partopsep=3pt\parskip=0pt
}
{\end{description}}

\newcommand{\jindent}[2]{
  \noindent\makebox[0pt][r]{{#1}\hspace*{\marginparsep}}
  \parbox[t]{0.95\linewidth}{#2}\par
}

%
\widowpenalty=1000
\clubpenalty=1000
%
%  XCSoar - Website Einfügen
\newcommand{\xcsoarwebsite}[1]{\url{http://www.xcsoar.org#1}}
%
% Define command to insert tip image
\newcommand{\tip}[0]{\marginlabel{\parbox{1.1cm}{\includegraphics[width=0.7cm]{figures/reminder.pdf}}}}

% Define command to insert gesture image
\newcommand{\gesture}[1]{\marginlabel{{\it#1{\phantom{aa}}}\parbox{1.1cm}{\hspace{-3mm}\includegraphics[width=0.8cm]{figures/gesture.pdf}}}}
%
% Define command to insert warning image
\newcommand{\warning}[0]{\marginlabel{\parbox{1.3cm}{\includegraphics[width=0.9cm]{figures/warning.pdf}}}}
%
% Define command to insert Achtung image
\newcommand{\achtung}[0]{\marginlabel{\parbox{1.3cm}{\includegraphics[width=2.5em]{figures/warning.pdf}}}}
%
% Define command to insert a flash image
\newcommand{\blitz}[0]{\marginlabel{\parbox{1.3cm}{\includegraphics[height=2.0em]{figures/reminder.pdf}}}}
%
% Define command to insert a stop
\newcommand{\halt}[0]{\marginlabel{\parbox{1.3cm}{\includegraphics[height=2.0em]{figures/warning.pdf}}}}
%
% Define command to reference a configuration item
\newcommand{\config}[1]{\marginlabel{\ref{conf:#1}\parbox{1.3cm}{\includegraphics[width=0.8cm]{figures/config.pdf}}}}
%
% Potentially overdue ``InfoBox'' style macro but sometimes used... so ..let it be alive ..
\newcommand{\InfoBox}[0]{{InfoBox}}
%
% Define command to put a menu label on the margin
\newcommand{\menulabel}[1]{\marginpar{\parbox{5.0cm}{\raggedright #1}}}
%
% Define command to draw a sketch on the margin
\newcommand{\sketch}[1]{\marginpar{\parbox{4.750cm}{\includegraphics[angle=0,width=0.9\linewidth,keepaspectratio='true']{#1}}}}
%
% Define command to draw a small sketch on the margin
\newcommand{\smallsketch}[1]{\marginpar{\includegraphics[angle=0,keepaspectratio='true']{#1}}}
%
% Enumerated todo's for the todonotes package
\newcounter{todocounter}
\newcommand{\todonum}[2][]{\stepcounter{todocounter}\todo[#1]{\thetodocounter: #2}}
%
% dies nette Makro bringt mir die aktuelle Version der bearbeiteten XCSoar-Verrion aufs Papier 
\newcommand{\version}{\begingroup\catcode`\_=\active\input{VERSION.txt}\endgroup}
%
%Anführungszeichen per Tastatur einfach so eingeben -> "
\shorthandoff{"}%


\title{Developer Manual}

\usepackage{calc}
\usepackage{fancyhdr}

\newcommand{\xcsoarheader}[1]{
\pagestyle{fancy}

% Add XCSoar User Manual title to the header
\fancyhead[L]{
\begin{tabular}{l}
\hspace*{-6cm} \em #1 \vspace*{2pt}
\end{tabular}
}

% Add page number to the footer (centered)
\fancyfoot{}
\fancyfoot[R]{\thepage}
}

% No line between content and header
\renewcommand{\headrulewidth}{0pt}

\fancypagestyle{plain}{
  % Clear fancy header and footer
  \fancyhf{}

  % Add page number to the footer (centered)
  \fancyfoot[R]{\thepage}

  % No line between content and header
  \renewcommand{\headrulewidth}{0pt}
}

\xcsoarheader{XCSoar Developer Manual}

\def\maketitle{%
  \null
  \thispagestyle{empty}%
  \begin{maxipage}
    \begin{center}
    \includegraphics[angle=0,width=\textwidth,keepaspectratio='true']{xcsoar-title.png}
    \end{center}
    \begin{center}
      \normalfont\huge\textsf{Glide~Computer and Navigation~System}\par
    \end{center}
    \vskip 1cm
    \begin{center}
      \normalfont\huge\textsf{\@title}\par
    \end{center}
    \vskip 1cm
  \end{maxipage}

  \vfill

  \begin{flushright}
    \large \strut {
      \sf
      \begin{tabular}{r}
      Manual version 1.6 \\
      \today \\
      For XCSoar version 6.0 \\
      \xcsoarwebsite \\
      \end{tabular} 
    } 
    \par
  \end{flushright}
  \par
  \vfil
  \vfil
  \null
  \cleardoublepage
}

\usepackage{listings}

\usepackage{tikz}
\usetikzlibrary{arrows,shapes,fit,decorations,decorations.pathmorphing,decorations.pathreplacing,calc}

\tikzstyle{thread}=[draw, fill=lightgray, text centered,
  minimum width=7em, minimum height=2.5em]

\begin{document}
\maketitle

\begingroup
\setlength{\cftsecnumwidth}{3em}
\tableofcontents
\endgroup


%%%%%%%%%%%%%%%%%%%%%%

\chapter*{Preface}

This manual applies to XCSoar version 7.0.  The authors reserve the
right to update this manual as enhancements are made throughout the
life of this product.

\section*{Warnings and precautions}

\warning IT IS THE USER'S RESPONSIBILITY TO USE THIS SOFTWARE PRUDENTLY. THIS
SOFTWARE IS INTENDED TO BE USED ONLY AS A NAVIGATION AID AND MUST NOT
BE USED FOR ANY PURPOSE REQUIRING PRECISE MEASUREMENT OF DIRECTION,
DISTANCE, LOCATION, OR TOPOGRAPHY. THIS SOFTWARE SHOULD NOT BE USED AS
AN AID TO DETERMINE GROUND PROXIMITY FOR AIRCRAFT NAVIGATION.
THIS SOFTWARE SHOULD NOT BE USED AS A TRAFFIC COLLISION AVOIDANCE SYSTEM.


\section*{Legal notices}

\subsection*{Software license agreement}

This software is released according to the GNU General Public License
Version~2.  See Appendix~\ref{cha:gnu-general-public} for the full
text of the agreement and warranty notice.

\subsection*{Limited liability}

In no event shall XCSoar, or its principals, shareholders, officers,
employees, affiliates, contractors, subsidiaries, or parent
organizations, be liable for any incidental, consequential, or
punitive damages whatsoever relating to the use of the Product.

\subsection*{Disclaimer}

This product, and all accompanying files, data and materials, are
distributed ``as is'' and with no warranties of any kind, whether
express or implied.  This product is used entirely at the risk of the
user.  Although great care has been taken to eliminate defects during
its development it is not claimed to be fault-free. No claims are made
regarding its correctness, reliability or fitness for any particular
purpose.  The XCSoar project developers and contributors shall not be
liable for errors contained herein or for incidental or consequential
damages, loss of data or personal injury in connection with
furnishing, performance, or use of this material.

%%%%%%%%%%%%%%%%%%%%%%%%%%%%%%%%%%%%%%%%%%%%%%%%%%%%%%%%%%%%%%%%

\chapter{Introduction}\label{cha:introduction}

\chapter{Compiling XCSoar}\label{cha:compiling}

The \texttt{make} command is used to launch the XCSoar build process.
You can learn more about the build system internals in chapter
\ref{cha:build}.

Most of this chapter describes how to build XCSoar on Linux, with
examples for Debian/Ubuntu.  A cross-compiler is used to build
binaries for other operating systems (for example Android and
Windows).

\section{Getting the Source Code}

The XCSoar source code is managed with
\href{http://git-scm.com/}{git}.  It can be downloaded with the
following command:

\begin{verbatim*}
git clone git://git.xcsoar.org/xcsoar/master/xcsoar.git
\end{verbatim*}

To update your repository, type:

\begin{verbatim*}
git pull
\end{verbatim*}

To update third-party libraries used by XCSoar (such as
\href{http://www.boost.org/}{Boost}), type:

\begin{verbatim*}
git submodule init
git submodule update
\end{verbatim*}

For more information, please read to the git documentation.

\section{Requirements}

The following is needed for all targets:

\begin{itemize}
\item GNU make
\item GNU compiler collection (\texttt{gcc}), version 4.8.1 or later
  or clang/LLVM 3.4 (with "make CLANG=y")
\item GNU gettext
\item \href{http://librsvg.sourceforge.net/)}{rsvg}
\item \href{http://www.imagemagick.org/}{ImageMagick 6.4}
\item \href{http://xmlsoft.org/XSLT/xsltproc2.html}{xsltproc}
\item Perl and XML::Parser
\item FFmpeg
\end{itemize}

The following command installs these on Debian:

\begin{verbatim*}
sudo apt-get install make \
  librsvg2-bin xsltproc \
  imagemagick gettext ffmpeg \
  git quilt
\end{verbatim*}

\section{Target-specific Build Instructions}

\subsection{Compiling for Linux/UNIX}

The following additional packages are needed to build for Linux and
similar operating systems:

\begin{itemize}
\item \href{http://www.zlib.net/}{zlib}
\item \href{http://curl.haxx.se/}{CURL}
\item \href{http://www.lua.org/}{Lua}
\item \href{http://www.libsdl.org/}{SDL}
\item \href{http://www.libsdl.org/projects/SDL\_ttf/}{SDL\_ttf}
\item \href{http://www.libpng.org/}{libpng}
\item \href{http://libjpeg.sourceforge.net/}{libjpeg}
\item OpenGL (Mesa)
\item to run XCSoar, you need one of the following fonts (Debian
  package): DejaVu (\texttt{fonts-dejavu}),
  Roboto (\texttt{fonts-roboto}),
  Droid (\texttt{fonts-droid}),
  Freefont (\texttt{fonts-freefont-ttf})
\end{itemize}

The following command installs these on Debian:

\begin{verbatim*}
sudo apt-get install make g++ \
  zlib1g-dev \
  libfreetype6-dev \
  libpng-dev libjpeg-dev \
  libtiff5-dev libgeotiff-dev \
  libcurl4-openssl-dev \
  liblua5.2-dev lua5.2-dev \
  libxml-parser-perl \
  libasound2-dev \
  librsvg2-bin xsltproc \
  imagemagick gettext \
  libegl1-mesa-dev \
  fonts-dejavu
\end{verbatim*}

To compile, run:

\begin{verbatim*}
make
\end{verbatim*}

You may specify one of the following targets with \texttt{TARGET=x}:

\begin{tabularx}{1.9\textwidth}{lX}

\texttt{UNIX} & regular build (the default setting) \\

\texttt{UNIX32} & generate 32 bit binary \\

\texttt{UNIX64} & generate 64 bit binary \\

\texttt{OPT} & alias for UNIX with optimisation and no debugging \\

\end{tabularx}

\subsection{Compiling for Android}

For Android, you need:

\begin{itemize}
\item \href{http://developer.android.com/sdk/}{Android SDK level 22}
\item \href{http://developer.android.com/sdk/ndk/}{Android NDK r13}
\item \href{http://www.vorbis.com/}{Ogg Vorbis}
\item \href{http://ant.apache.org/}{Apache Ant}
\item {Java Native tools 
\begin{verbatim*}
sudo apt-get install gcj-native-helper ant vorbis-tools
\end{verbatim*}}
\end{itemize}

The \texttt{Makefile} assumes that the Android SDK is installed in
\verb|~/opt/android-sdk-linux| and the NDK is installed in
\verb|~/opt/android-ndk-r13|.  You can use the options
\verb|ANDROID_SDK| and \verb|ANDROID_NDK| to override these paths.

After installing Java you will have to install at least one Android Platform SDK using the Android SDK Manager:

\begin{verbatim*}
~/opt/android-sdk-linux/tools/android
\end{verbatim*}

Load/update the IOIO source code:

\begin{verbatim*}
git submodule init
git submodule update
\end{verbatim*}

To compile, run:

\begin{verbatim*}
make TARGET=ANDROID
\end{verbatim*}

Use one of the following targets:

\begin{tabularx}{1.9\textwidth}{lX}

\texttt{ANDROID} & for ARM CPUs (same as \texttt{ANDROID7}) \\

\texttt{ANDROID7} & for ARMv7 CPUs \\

\texttt{ANDROID7NEON} & with
\href{http://www.arm.com/products/processors/technologies/neon.php}{NEON}
extension \\

\texttt{ANDROID86} & for x86 CPUs \\

\texttt{ANDROIDMIPS} & for MIPS CPUs \\

\texttt{ANDROIDFAT} & "fat" package for all supported CPUs \\

\end{tabularx}

It is possible that problems occur on 64bit OS (which is now required for Android SDK) when you try to run the cross-compiler toolchain which contains 32 bit components. In some distributions this can be fixed by using the \emph{ia32-libs}, newer distributions (Ubuntu 16.04 for example) might require a different approach (e.g. \emph{lib32z1}).

\subsection{Compiling for Windows}

To cross-compile to (desktop) Windows, you need the mingw-w64 version
of gcc:

 http://mingw-w64.sourceforge.net/

To compile, run one of the following:

\begin{verbatim*}
make TARGET=PC
\end{verbatim*}

Use one of the following targets:

\begin{tabularx}{1.9\textwidth}{lX}

\texttt{PC} & 32 bit Windows (i686) \\

\texttt{WIN64} & Windows x64 (amd64 / x86-64) \\

\texttt{CYGWIN} & Windows build with Cygwin (experimental) \\

\end{tabularx}

\subsection{Compiling for iOS and OS X}

On OS X, the following tools are required:
\begin{itemize}
\item png2icns from \href{http://icns.sourceforge.net}{libicns} to build for
  OS X
\item \href{https://alioth.debian.org/projects/dpkg}{dpkg} to build the iOS
  Cydia package
\item \href{http://cdrecord.org/private/cdrecord.html}{mkisofs} to build the
  OS X DMG package
\end{itemize}

To compile for iOS / AArch64, run:

\begin{verbatim*}
make TARGET=IOS64 cydia-deb
\end{verbatim*}

To compile for iOS / ARMv7, run:

\begin{verbatim*}
make TARGET=IOS32 cydia-deb
\end{verbatim*}

To compile for OS X / x86\_64, run:

\begin{verbatim*}
make TARGET=OSX64 dmg
\end{verbatim*}

\subsection{Compiling for the Raspberry Pi}

You need an ARM toolchain.  For example, you can use the Debian
package \verb|g++-5-arm-linux-gnueabihf|:

\begin{verbatim*}
make TARGET=PI TCSUFFIX=-5
\end{verbatim*}

To optimize for the Raspberry Pi 2 (which has an ARMv7 with NEON
instead of an ARMv6):

\begin{verbatim*}
make TARGET=PI2 TCSUFFIX=-5
\end{verbatim*}

These targets are only used for cross-compiling on a (desktop)
computer.
If you compile on the Raspberry Pi, the default target will
auto-detect the Pi.

\subsection{Compiling for the Cubieboard}

To compile, run:

\begin{verbatim*}
make TARGET=CUBIE
\end{verbatim*}

This target is only used for cross-compiling on a (desktop) computer.
If you compile on the Cubieboard, the default target will auto-detect
the Cubieboard.

\subsection{Compiling for Kobo E-book Readers}

You need an ARM toolchain.  For example, you can use the Debian
package \verb|g++-5-arm-linux-gnueabihf|.

To compile XCSoar, run:

\begin{verbatim*}
make TARGET=KOBO TCSUFFIX=-5
\end{verbatim*}

To build the kobo install file \texttt{KoboRoot.tgz}, you need the
following Debian packages:

\begin{verbatim*}
sudo apt-get install fakeroot ttf-bitstream-vera
\end{verbatim*}

Then compile using this command:

\begin{verbatim*}
make TARGET=KOBO TCSUFFIX=-5 output/KOBO/KoboRoot.tgz
\end{verbatim*}

For this, you need the Debian package \verb|libc6-armhf-cross|.

\subsubsection{Building USB-OTG Kobo Kernel}

To build a USB-OTG capable kernel for the Kobo, clone the git
repository:

\begin{verbatim*}
git clone git://git.xcsoar.org/xcsoar/max/linux.git
\end{verbatim*}

Check out the correct branch. For the Kobo Mini, this is the "kobo" branch,
for the Kobo Glo HD, the branch is called "kobo-glohd", and for the Kobo Aura 2,
use the branch "kobo-aura2".

\begin{verbatim*}
git checkout kobo
\end{verbatim*}

Configure the kernel using the configuration files from the \texttt{kobo/kernel}
directory in XCSoar's \texttt{git} repository. For the Kobo Mini, install a
\href{http://openlinux.amlogic.com:8000/download/ARM/gnutools/arm-2010q1-202-arm-none-linux-gnueabi-i686-pc-linux-gnu.tar.bz2}{gcc
  4.4 cross compiler}, for example in \texttt{/opt}. For the Kobo Glo HD and
Aura 2, install a
\href{https://launchpad.net/gcc-arm-embedded/4.6/4.6-2012-q4-update/+download/gcc-arm-none-eabi-4_6-2012q4-20121016.tar.bz2}{gcc
  4.6 cross compiler}

To compile a kernel image for the Kobo Mini, type:

\begin{verbatim*}
make \
  CROSS_COMPILE=/opt/arm-2010q1/bin/arm-none-linux-gnueabi- \
  ARCH=arm uImage
\end{verbatim*}

To compile a kernel image for the Kobo Glo HD, type:

\begin{verbatim*}
make \
  CROSS_COMPILE=/opt/gcc-arm-none-eabi-4_6-2012q4/bin/arm-none-eabi- \
  ARCH=arm uImage
\end{verbatim*}

Copy \texttt{uImage} to the Kobo.  Kernel images can be installed with
the following command:

\begin{verbatim*}
dd if=/path/to/uImage of=/dev/mmcblk0 bs=512 seek=2048
\end{verbatim*}

Note that XCSoar's \texttt{rcS} script may overwrite the kernel image
automatically under certain conditions.  To use a new kernel
permanently, install it in \texttt{/opt/xcsoar/lib/kernel}.  Read the
file \texttt{kobo/rcS} to find out more about this.

To include kernel images in \texttt{KoboRoot.tgz}, copy
\texttt{uImage.otg}, \texttt{uImage.kobo}, \texttt{uImage.glohd.otg},
\texttt{uImage.glohd}, \texttt{uImage.aura2} and \texttt{uImage.aura2.otg} to
\texttt{/opt/kobo/kernel}.

\subsection{Editing the Manuals}

The XCSoar documententation, including the Developer Manual that you are 
reading right now, is written using the TeX markup language. You can edit
the source files with any text editor, although a specific TeX editor (e.g. LateXila) 
makes it easier. 

Source files are located in the en, fr, de, pl subdirectories of the doc/manual directory. 
The Developer manual is in the doc/manual/en directory. The generated files are put into
the output/manual directory.

To generate the PDF manuals, you need the TexLive package, plus some European languages.

The following command installs these on Debian:

\begin{verbatim*}
sudo apt-get install texlive \
  texlive-latex-extra \
  texlive-luatex \
  texlive-lang-french \
  texlive-lang-polish \
  texlive-lang-german \
  liblocale-po-perl
\end{verbatim*}

The documentation is distributed as PDF files. Generating the PDFs from the 
TeX files is done by typing:

\begin{verbatim*}
make manual
\end{verbatim*}

A lot of warnings are generated... this is normal. Check for the presence of PDF files 
to ensure that the generation process was successful.
\section{Options}

\subsection{Parallel Build}

Most contemporary computers have multiple CPU cores.  To take
advantage of these, use the \texttt{make -j} option:

\begin{verbatim*}
make -j12
\end{verbatim*}

This command launches 12 compiler processes at the same time.

Rule of thumb: choose a number that is slightly larger than the number
of CPU cores in your computer.  12 is a good choice for a computer
with 8 CPU cores.

\subsection{Optimised Build}

By default, debugging is enabled and compiler optimisations are
disabled.  The resulting binaries are very slow.  During development,
that is helpful, because it catches more bugs.

To produce optimised binaries, use the option \texttt{DEBUG}:

\begin{verbatim*}
make DEBUG=n
\end{verbatim*}

Be sure to clean the output directory before you change the
\texttt{DEBUG} setting, because debug and non-debug output files are
not compatible.

The convenience target \texttt{OPT} is a shortcut for:

\begin{verbatim*}
TARGET=UNIX DEBUG=n TARGET_OUTPUT_DIR=output/OPT
\end{verbatim*}

It allows building both debug and non-debug incrementally, because two
different output directories are used.

\subsection{Compiling with ccache}

To speed up the compilation of XCSoar we can use \texttt{ccache} to cache the 
object files for us. All we have to do is install ccache and 
add 
\texttt{USE\_CCACHE=y} 
to the make command line:

\begin{verbatim*}
sudo apt-get install ccache
make TARGET=UNIX USE_CCACHE=y
\end{verbatim*}


\chapter{Policy}\label{cha:policy}

\section{Git Work Flow}

\subsection{Version Numbering}

Each release of XCSoar is denoted by a unique version number. A version
number consists three sequential numerical fields separated a period and
prefixed by a 'v'. The significance of the three fields are
<major>.<minor>.<patch>. So, for instance, the version number 'v7.1.3'
indicates XCSoar major version 7, minor version 1, and patch 3.

Release versions are determined and maintained by the release manager.

\subsection{Git Repository Enduring Branches}

At any time the XCSoar contains two enduring branches. The principle
enduring branch is 'master'. The lifetime of this branch is unbounded.
The second enduring branch is the current minor version branch. The
lifetime of this branch is the lifetime of the current minor version.
The name of the minor version branch is in the form 'v<major>.<minor>.x'
So, for instance, the current minor branch may be called 'v7.1.x'.

With the commencement of a new minor version the release manager will
create the required branch. The old minor branch will be retired.
All XCSoar releases are made from the current minor branch.

The master branch serves as the development branch for the next minor
release (which may also be the first minor of the next major release).

Developers should commit changes for the next minor release to the master
branch. Bug fixes of the current minor release should be committed to
the current minor branch, in preparation for the next patch release.

\section{Writing Patches}

There are two methods to summit patches to the upstream XCSoar repository.

\subsection{GitHub}
The XCSoar repository is currently hosted at
\href{https://github.com/XCSoar/XCSoar}{XCSoar GitHub repository}.
GitHub provides a method for submitting patches using their Web interface.
The easiest way to use this method is to make another XCSoar repository by
 'forking' from 'XCSoar/XCSoar'. Make sure that your local repository
is up to date with the upstream XCSoar/XCSoar repository. Then make a 
feature branch from the master or current minor branch as appropriate for
the change you are making and commit the changes to that branch. To make
the upstream merge easier it is best to re base this feature branch with
the appropriate upstream branch from time to time.

When ready to submit the feature branch to the upstream push your local
repository to your GitHub XCSoar repository and use the GitHub Web interface
to make a 'pull request' of your feature branch to XCSoar/SCSoar.

\subsection{Developers' Mail List}
The second method is to submit patches or \texttt{git pull} requests to
the developer mailing list (\texttt{xcsoar-devel@lists.sourceforge.net}).

Patch files can be generated by running
\begin{verbatim*}
git diff > patch
\end{verbatim*}

\subsection{Basic Patch Requirements}
A patch should be self-explanatory, it needs a good description.  The
subject line specifies the subsystem/library name and a brief
description of what is changed, followed by an empty line.  Then write
a longer description if needed, and explain why this change is needed.

Each patch must compile and must not introduce a regression (as far as
we know at the time).

Each patch must be self-contained and should only change one thing.
Split larger patches into smaller pieces.  Don't refactor and
add/modify/remove features in the same patch.

Don't rewrite code unless you need to.  Migrate incrementally to a new
concept.  Keep patches small and easy to understand.

\section{Code Style}

79 columns, reasonable exceptions allowed.  Indent 2 spaces, no tabs.
No indent for \verb|namespace| blocks (a compromise to avoid excessive
indentation).

\emph{Comments:} write enough code comments (in English).  All
workarounds must be documented.  Everybody must be able to understand
your code, even when you're gone.  Don't abuse multiple single-line
comments (``//'') to write mult-line comments.

\emph{API documentation:} non-trivial functions should be documented
in a doxygen comment.

\emph{Names:} class/function names in \texttt{CamelCase} (not
\texttt{camelCase}); attributes/variables lower case, separated with
underscore (e.g. \texttt{foo\_bar}); constants (including
\texttt{enum} values) all upper case (e.g. \texttt{FOO\_BAR}).

Exception: when a foreign API is being mimicked (e.g. STL containers),
we adopt its naming conventions.

\emph{Files:} \texttt{*.cpp} and \texttt{*.hpp} for C++.  Files should
be named after the main class which is provided.  Each class should
have a separate source file and a separate header.  UNIX text format.

Be \texttt{const}-correct.  Use \verb|constexpr| instead of
\verb|const| whenever possible.

Use \verb|static| whenever possible.  Functions and global variables
that are only used in one source file should not be exported.  Methods
that do not use any instance method/variable should be static to avoid
the overhead of passing the implicit \verb|this| parameter.

Make methods \verb|virtual| only after careful consideration.  A
destructor should only be virtual if necessary.  All overrides must
use the \verb|override| keyword.  Use \verb|final| often.

Compile with \texttt{WERROR=y} and fix all warnings.

Don't write large functions.  Split them up when they become too
large.

Avoid dynamic allocation.  Dynamic allocation means overhead, more
locking and heap fragmentation.  Use \texttt{StaticArray} and
\texttt{StaticString} if possible.

Asterisks belong to the variable name, not to the type name.  Consider
``\texttt{Foo* a, b}''.  ``\texttt{Foo *a, b}'' or ``\texttt{Foo *a,
  *b}'' is easier to understand.

Some sample code to demonstrate our code style:

\begin{verbatim}
/**
 * API documentation for this class.
 */
struct TheStruct {
  unsigned an_attribute;
  bool second_attribute;

  TheStruct();

  /**
   * API documentation for this method.
   *
   * @param foo documentation for this parameter
   * @return documentation for the return value
   */
  bool TheMethod(int foo);
};

TheStruct::TheStruct()
  :an_attribute(0),
   second_attribute(true)
{
}

static bool
FooBar(int a_parameter, unsigned another_parameter,
       const TheStruct *next_row)
{
  switch (a_parameter) {
  case 0:
    break;
  }

  if (a_parameter == 2 && another_parameter == 3 &&
      next_row != NULL)
    return true;

  return a_parameter == 42;
}
\end{verbatim}

\section{C++}

XCSoar is written in C++17.

XCSoar's standard compilers are \texttt{gcc} (at least version 6) and
\texttt{clang} (at least version 4.0).

Avoid preprocessor macros, because they are obscure, error prone, not
type-safe, hard to read and hard to debug.  Use \texttt{inline}
functions and \texttt{constexpr} variables instead.

\subsection{Other rules}

In a class declaration, attributes come first, then
constructor/destructor, and finally the methods.  Having all
attributes in one place gives a good overview of the nature of a
class.

Avoid expensive and bloated STL containers if there are cheaper
solutions (e.g. \texttt{StaticArray}, \texttt{StaticString} if the
maximum size is predictable).

Avoid template hell.  Keep templates readable.  Keep in mind that
excessive template use may bloat the binary.

\section{Graphical User Interface}
\subsection{Letter Cases}

Following the guidline should prevent the GUI from mixtures of ``ON'' and
``On'' text elements, and lead to a systematic GUI text presentation.
The goal is to recognize GUI text fast and reliable.

\begin{description}
  \item[Captions] Captions (button captions, windows titles) to use
        capitalization. E.g. ,``Pan On'', ``The Display Of ...''.

  \item[Abbreviations] Generally known abbreviation use upper case like ``MC'', ``ETA'',
        ``V''; or they can use CamelCase, especially when using synthetic
        words  like ``GoTo'', ``InfoBox''.
        Abbreviated words by simply cutting the end of the word needs a
        dot, e.g. ``Max.\ temp.''

  \item[Plain text] Longer help texts are to write like prose: ``This is the help
        page for ...''.

  \item[Labels] Label text has the least systematic constraints:
\begin{itemize}
    \item Captions for text (input) fields, e.g. ``Wing loading''
    \item Info text on widgets. E.g. ``No data'' on an empty
          analysis page. 
    \item Label text for radio or check boxes. 
    \item Selections on Combo-boxes, selectors, Pull-down menus. 
\end{itemize}
  All those should go like prose, whereas exceptions might be
  meaningful.
  
  \item[Gauge caption] Also the appearance of the gauge caption should be covered with
        that. They are currently mapped to upper case all over. I think
        the most readable also here is a CamelCase approach. E.g.\ to
        distinct ``WP Dist'', ``WP AltD'', and ``WP AltR''. Another good
        example would be MACCREADY, which should be MacCready, or just
        MC.

  \item[Units] Units have their own specific appearance. A profound paper is
        http://physics.nist.gov/cuu/pdf/checklist.pdf we could just
        refer to.
\end{description}


\chapter{Architecture}

This chapter describes XCSoar's internal code architecture.

\section{Source Organisation}

XCSoar's source code is stored in the \texttt{src} directory.  This
section tries to give a rough overview where you can find what.

\begin{itemize}

\item \texttt{Util/}: generic C++ utilities that do not depend on
  external libraries, such as data structures, string operations

\item \texttt{Math/}: math data types (fixed-point math, angles) and
  generic formulas

\item \texttt{Geo/}: geographic data structures and formulas

\item \texttt{Formatter/}: code that formats internal values to
  strings

\item \texttt{Units/}: conversion from SI units (``System'' units) to
  configured user units

\item \texttt{NMEA/}: data structures for values parsed from NMEA

\item \texttt{Profile/}: user profiles, loading from and saving to

\item \texttt{IGC/}: support for the IGC file format

\item \texttt{Logger/}: all loggers (NMEA, IGC, flights)

\item \texttt{Thread/}: multi-threading support (OS specific)

\item \texttt{Screen/}: base library for the graphical user interface

\item \texttt{Renderer/}: various graphical renderers, for map and
  analysis

\item \texttt{MapWindow/}: the map

\item \texttt{Form/}: modal dialogs and their controls (based on the
  screen library)

\item \texttt{Dialogs/}: modal dialogs implementations (based on the
  form library)

\item \texttt{Net/}: networking code (OS specific)

\item \texttt{Operation/}: generic code to support cancellable
  long-running operations

\item \texttt{Android/}: code specific to Android (the native part
  only; Java code is in \texttt{android/src/}

\item \texttt{Engine/PathSolvers/}: an implementation of Dijkstra's
  path finding algorithm, for task and contest optimisation

\item \texttt{Engine/Airspace/}: airspace data structures and airspace
  warnings

\item \texttt{Engine/Waypoint/}: waypoint data structures

\item \texttt{Engine/GlideSolvers/}: a MacCready implementation

\item \texttt{Engine/Task/}: task data structures and calculations

\item \texttt{Engine/Contest/}: contest optimisation

\item \texttt{Engine/Route/}: the route planner (airspace and terrain)

\end{itemize}

\section{Threads and Locking}

\subsection{Threads}

XCSoar runs on multiple threads, to make the UI responsive but still
allow expensive background calculations.

This is how it looks like on Windows and Linux/SDL (software
rendering):

\begin{tikzpicture}
  \node(ui)[thread]{UI thread};

  \path(ui.east)+(4,0) node(draw)[thread]{DrawThread};
  \draw[->] (ui.10) -- (draw.170)
  node[midway, above, sloped, font=\tiny]{redraw (Pan)};
  \draw[->] (draw.190) -- (ui.350)
  node[midway, below, sloped, font=\tiny]{BufferCanvas};

  \path(draw.south)+(0,-2) node(calc)[thread]{CalcThread};
  \draw[->] (calc.north) -- (draw.south)
  node[midway, below, sloped, font=\tiny]{results};

  \path(calc.south)+(0,-2) node(merge)[thread]{MergeThread};
  \draw[->] (merge.north) -- (calc.south)
  node[midway, below, sloped, font=\tiny]{results};

  \path(merge.south)+(-3,-2) node(device1)[thread]{Device};
  \draw[->] (device1.north) -- (merge.250)
  node[midway, above, sloped, font=\tiny]{sensor data};

  \path(merge.south)+(3,-2) node(device2)[thread]{Device 2};
  \draw[->] (device2.north) -- (merge.290)
  node[midway, above, sloped, font=\tiny]{sensor data};

  \path(ui.south)+(0,-3) node(io)[thread]{I/O thread};
  \draw[->] (io.north) -- (ui.south)
  node[midway, above, sloped, font=\tiny]{data};
  \draw[->] (io.east) -- (merge.west)
  node[midway, above, sloped, font=\tiny]{sensor data};
\end{tikzpicture}

The UI thread is the main thread.  It starts the other threads and is
responsible for the UI event loop.  No other thread is allowed to
manipulate windows.  The UI thread has a timer which does regular
house keeping twice per second (\texttt{Pro\-cess\-Ti\-mer.cpp}).

The calculation thread (\texttt{Calcu\-la\-tion\-Thread.cpp},
\texttt{Glide\-Com\-pu\-ter*.cpp}) does all the expensive calculations
in background.  It gets data from the devices (through
\texttt{Merge\-Thread}) and forwards it together with calculation
results to the drawing thread and the main thread.

Each device has its own thread (\texttt{Serial\-Port.cpp}).  This is
needed because Windows CE does not support asynchronous COMM port I/O.
The thread is stopped during task declaration (which happens in the UI
thread).

When new data arrives on the serial port, the \texttt{Merge\-Thread}
gets notified, which will merge all sensor values into one data
structure.  It will then run cheap calculations, and forwards
everything to the \texttt{Calcu\-la\-tion\-Thread}.

With OpenGL, the map is rendered live without a buffer.  There is no
DrawThread.

On Android, the UI thread is not the main thread - the main thread is
implemented in Java, managed by Android itself.  The UI thread listens
for events which the Java part drops into the event queue
(\texttt{NativeView.java} and others).  The internal GPS does not need
a thread, it is implemented with Java callbacks.  For Bluetooth I/O,
there are two threads implemented in Java (\texttt{InputThread.java}
and \texttt{OutputThread.java}, managed by
\texttt{BluetoothHelper.java}).

\subsection{Locking}

Some data structures are rarely modified.  There is no lock for them.
For a modifications, all threads must be suspended.  Example:
waypoints, airspaces.

Other data structures are modified so often that correct locking would
be too much overhead.  Each thread and each instance has its own
copy.  The lock needs to be obtained only for making the private
copy.  The private copy can be used without locking.  Example:
\texttt{NMEA\_INFO}, \texttt{DERIVED\_INFO}.

There are objects which are too expensive to copy.  Normal locking
applies to them.  We have a template class called \texttt{Guard} to
enforce proper read/write locking.  Example: the task.


\section{Accessing Sensor Data}

Much of XCSoar deals with obtaining sensor data and visualising it.

Suppose you want to write a dialog that needs the current GPS
location, where do you get it?  The short and simple answer is: from
\verb|CommonInterface::Basic()| (the \texttt{InterfaceBlackboard}).
Example:

\begin{verbatim}
#include "Interface.hpp"

...
  const auto &basic = CommonInterface::Basic();
  if (basic.location_available)
    current_location = basic.location;
\end{verbatim}

This is true for the main thread (aka the ``user interface thread'').
Other threads must not use the \texttt{Interface.hpp} library, because
the \verb|InterfaceBlackboard| is not protected in any way.  It
contains copies of various data structures just for the main thread.

This is how sensor data moves inside XCSoar:

\begin{tikzpicture}
  \node(merge)[thread,align=center]{
    MergeThread\\
    BasicComputer\\
    DeviceBlackboard
  };

  \path(merge.north)+(-2,2) node(device1)[thread]{Device 1};
  \draw[->] (device1.south) -- (merge.north)
  node[midway, below, sloped, font=\tiny]{NMEAInfo};

  \path(merge.north)+(2,2) node(device2)[thread]{Device 2};
  \draw[->] (device2.south) -- (merge.north)
  node[midway, below, sloped, font=\tiny]{NMEAInfo};

  \path(merge.south)+(0,-2) node(calc)[thread,align=center]{
    CalcThread\\
    GlideComputer\\
    GlideComputerBlackboard
  };
  \draw[->] (merge.south) -- (calc.north)
  node[midway, below, sloped, font=\tiny]{MoreData};

  \path(calc.south)+(0,-2) node(ui)[thread,align=center]{
    UI thread\\
    InterfaceBlackboard\\
    BlackboardListener
  };
  \draw[->] (calc.south) -- (ui.north)
  node[midway, below, sloped, font=\tiny]{DerivedInfo};

  \path(ui.east)+(3.5,0) node(draw)[thread,align=center]{
    DrawThread\\
    MapWindow\\
    MapWindowBlackboard
  };
  \draw[->] (ui.east) -- (draw.west)
  node[midway, below, sloped, font=\tiny]{DerivedInfo};
\end{tikzpicture}

The device driver parses input received from its device into its own
\texttt{NMEAInfo} instance inside \texttt{DeviceBlackboard}
(i.e. \verb|per_device_data|).  Then it wakes up the
\texttt{MergeThread} to merge the new data into the central
\texttt{NMEAInfo} instance.  The \texttt{MergeThread} hosts the
\texttt{BasicComputer} which attempts to calculate missing data (for
example, derives vario from GPS altitude).

The \texttt{CalculationThread} wakes up and receives the
\texttt{MoreData} object from \texttt{DeviceBlackboard}.  Here,
expensive calculations are performed (\texttt{GlideComputer}: task
engine, airspace warnings, ...), resulting in a \texttt{DerivedInfo}
object.  The \texttt{CalculationThread} runs no more than twice per
second.

Finally, the UI thread wakes up and receives \texttt{MoreData} and
\texttt{DerivedInfo} via \texttt{DeviceBlackboard}.  This updates
InfoBoxes and other UI elements.  On Windows, the map is drawn in a
separate thread, so there's another layer.

Let's get back to the question: where do I get sensor data?  That
depends on who you are:

\begin{itemize}

\item you are the user interface: (InfoBoxes, dialogs, any Window
  callback): \texttt{InterfaceBlackboard} (see above).  To get
  notified on changes, register a \texttt{BlackboardListener} (and
  don't forget to unregister it).

\item you are the MapWindow: depends!  If you're being called from
  \texttt{OnPaintBuffer} (i.e.\ inside the \texttt{DrawThread}), you
  must use the \texttt{MapWindowBlackboard}, all others must use the
  \texttt{InterfaceBlackboard}.

\item you are a ``computer'' library: you will get the values as a
  parameter.  Don't try to use the \texttt{GlideComputerBlackboard}
  directly.

\item you are a device driver: implement the method
  \texttt{OnSensorUpdate} or \texttt{OnCalculatedUpdate} if you need
  to know values from other devices or calculation results.

\item everybody else may use the \texttt{DeviceBlackboard}, but be
  sure to lock it while using its data.

\end{itemize}


\chapter{The build system}\label{cha:build}

A big plain \texttt{Makefile} is used to control the XCSoar build.
GNU extensions are allowed.

This chapter describes the internals of our build system; for
instructions on compiling XCSoar, see chapter \ref{cha:compiling}.

\section{Linker parameters}

The following variables (or variable suffixes) appear in the
\texttt{Makefile} (conforming to \texttt{automake} conventions):

\begin{description}
\item[\texttt{LDFLAGS}] Linker flags, such as \texttt{-static} or
  \texttt{-Wl,...}, but not \texttt{-l}.
\item[\texttt{LDLIBS}] All \texttt{-l} flags, e.g. \texttt{-lGL}.
\item[\texttt{LDADD}] Path names of static libraries,
  e.g. \texttt{/usr/lib/libz.a}.
\end{description}

Search directories (\texttt{-L}) are technically linker ``flags'', but
they are allowed in \texttt{LDLIBS}, too.


\chapter{Developing}

\section{Debugging XCSoar}

The XCSoar source repository contains a module for the GNU debugger
(\texttt{gdb}).  It contains pretty-printers for various XCSoar types,
including \texttt{Angle}, \texttt{GeoPoint} and others.  These are
helpful when you print values in the debugger.  To use it, start the
debugging session and load the module:

\begin{verbatim}
$ gdb -ex "source tools/gdb.py" output/UNIX/bin/xcsoar
(gdb) run
\end{verbatim}

The module will automatically convert fixed-point to floating point,
radian angles to degrees and more.  You can now do fancy stuff like:

\begin{verbatim}
(gdb) p basic.location
$1 = GeoPoint(7.93911242887 51.1470221074)
(gdb) p basic.date_time_utc
$2 = DateTime(2012/12/23 21:41:57)
(gdb) p basic.track
$3 = 55.2254197961
(gdb) p basic.external_wind
$4 = GeoVector::ZERO
(gdb) p current_leg.vector_remaining
$5 = GeoVector(267.899420345 107957.109724)
\end{verbatim}



%%%%%%%%%%%%%%%%%%%%%%%%%%%%%%%%%%%%%%%%%%%%%%%%%%%%%%%%%%%%%%%
\chapter{User interface guidelines}

\section{General}

\begin{itemize}
\item Minimise the number of colours, and re-use colour groups already defined.
\item Too much use of colour where it is not required serves only to reduce 
 the effectiveness of bright colours for important items.
\item High colour saturation elements should be reserved for high importance items
\item High contrast against background should be reserved for high importance items
\item Attempt to adopt colours that are intuitive based the function of the item
\item Minimise the clutter where possible --- readibility is essential for use 
 in flight
\item Use colours defined in \verb|Graphics| according to functional name, not 
 their actual colour.
\item Try to maintain consistent use of colours in all uses of that function, 
 such as dialogue graphics as well as map overlays and infoboxes.
\item Text should always be monochrome.
\end{itemize}

Use aviation conventions or adopt best aviation human factors
standards where possible, in particular:
\begin{itemize}
\item ICAO Internation Standards and Recommended Practices, Annex 4 to the 
 Convention on International Civil Aviation (Aeronautical Charts).
\item NASA Colour Usage recommendations and design guidelines: \verb|http://colorusage.arc.nasa.gov/|
\item DOT/FAA/AR-03/67 {\em Human Factors Considerations in the Design and 
 Evaluation of Electronic Flight Bags (EFBs)} \verb|http://www.volpe.dot.gov/hf/aviation/efb/docs/efb_version2.pdf|
\item FAA Human Factors Design Standards \verb|http://hf.tc.faa.gov/hfds/|.
\item DOT/FAA/AM-01/17 {\em Human Factors Design Guidelines for Multifunction Displays}
\end{itemize}

Check for performance with respect to colour blindness.
This site has a useful tool that can be used to convert screenshots
to how they would look to a person with common color blindness:
\verb|http://www.etre.com/tools/colourcheck/|.

{\bf For safety purposes, avoid use of elements that may encourage or require 
the user to stare at the screen continuously.}

{\bf For safety purposes, avoid user controls that have significant risk of 
producing unsafe results if misconfigured by the pilot.}

\subsection{General colour conventions}
Colour conventions generally in use throughout the program:
\begin{itemize}
\item Red for indicator of warning
\item Orange for indicator of caution
\item Green for positive indicator of safety
\item Blue for neutral indicator of safety
\end{itemize}

\subsection{Displayed data}
\begin{itemize}
\item Where data is invalid, indicate this by not presenting the data or
  showing dashes.
\item Present data in user-defined units.
\item Display numerical data with significant digits appropriate to the accuracy of the
  calculations, or its functional use by the pilot, whichever is lower.
\end{itemize}

\section{Dialogs and menu buttons}

\subsection{Colors}
Colour conventions in use are:
\begin{itemize}
\item Grey for buttons
\item Buttons and other widgets rendered with an evenly shaded border
\item Yellow for clicked items
\item Light blue for the key focused item
\item Medium blue for dialogue title bar
\item Text is black if the item is enabled
\item Text is greyed out (but still visible) if the item is disabled
\end{itemize}

\subsection{dialogue types and navigation buttons}
There are four types of dialogs in XCSoar, and the navigation 
buttons for each are different.  
Navigation buttons are the Close, OK, Cancel and Select buttons.
\begin{itemize}
\item Dialogs that modify and save data when the dialogue closes.

These shall usually have a Close button (no Cancel) and may have context specific 
function buttons
\item Dialogs that modify data where Cancel would be important for the user.

These shall have OK and Cancel buttons.  This may include dialogs with 
children dialogs where hitting Cancel from the parent dialogue cancels 
all the changes made in the children dialogs

\item Dialogs that have a list of values, one of which can be selected to 
 return to the parent dialogue.

These shall have Select and Cancel buttons
\item Dialogs that display information that cannot be modified.

These shall have a Close button
\end{itemize}

\subsection{dialogue button placement and size}
\begin{itemize}
\item The Close and Cancel buttons will never appear in the same dialogue and 
 are always located in the same place.  This location will be:

For portrait: lower right

For landscape: lower left
\item The Select button will be accompanied with a Cancel button.  The locations will be:

For portrait: Select in lower left, Cancel in lower right

For landscape: Cancel in lower left, Select immediately above it
\item Buttons will be 35 (scaled) pixels high
\item Buttons will be flush with the bottom of the screen and with the sides of 
 the screen and against each other (no margins)
\item In portrait, buttons will be 33% or 50% or 100% width of the screen
\item In landscape, buttons will be 65 to 80 (scaled) pixels wide, as wide as the 
 frame permits.  They will generally be a vertical row of buttons flush left of the screen
\item If text won't fit on a button, the buttons can be made larger consistently 
 for a screen, but this should be the exception because if it must contain that 
 much text consider using a different type of control.
\item Exceptions to all the dialogue concepts above are encouraged, but should be 
 mocked up and reviewed with the development community prior to implementing and 
 possibly documenting in the developers guide.
\end{itemize}

\subsection{Usability}
\begin{itemize}
\item Minimum size of buttons should be X by Y mm
\item Ensure all dialogs are navigable using cursor keys only
\item Ensure the focussed item is clearly identified.  The rectangle
  of the widget on the canvas may be drawn using the \verb|fill_focus| method
  of \verb|Canvas|.
\end{itemize}

\section{Main graphics}

\subsection{Colors}
Colour conventions in use, in order of priority, are:
\begin{itemize}
\item Aircraft black and white, for neutrality but clear identification
\item Traffic (FLARM) use alarm green, orange, and red.
\item Lift is vibrant green, sink is copper orange.
\item Aircraft navigation (route, best cruise track) is (ICAO) dark purple-blue
\item Task navigation lines and areas are (ICAO) magenta.
\item Updraft sources and other updraft derived data is sky blue.
\end{itemize}

(Todo) airspace alert colours

Map culture (topography) and terrain rendering should conform to ICAO
Annex 4 where appropriate.  Note that some modifications are
reasonable for electronic use given that Annex 4 deals with paper
charts.  Nevertheless, the colour conventions are useful to adopt as they are
likely to be intuitive and are designed for aviation use.

\subsection{Pen styles}

\begin{itemize}
\item Map culture should be rendered with a thin pen
\item Thicker pens used for important (e.g.\ task, navigational, airspace) lines
\item Dashed lines are used to increase perceptual priority
\end{itemize}

\subsection{Map overlays}
Elements on the map that are not part of the map layer, such as additional
informational widgets (final glide bar, wind, north arrow) should be rendered
so as to help those elements be visually separated from the map:

\begin{itemize}
\item Generally adopt higher contrast (higher colour saturation or darker shade) 
 than the background map layer elements.

\item For elements covering an area (non line), draw the entire element or a border
with a luminosity contrasting pen, of width \verb|IBLSCALE(1)|.

\item Consider whether the widget is required in all flying states and display modes.
if it does not serve a direct functional purpose in some states/modes, do not
render it.

\item Avoid locating widgets at the aircraft symbol (ownship symbol).
It is important to keep this area clear so the aircraft symbol can be easily found.

\end{itemize}

Elements that may be rendered over each other should be organised in order of
priority, particularly with alert warning items above caution items above non-alert items.

\section{Terminology}
\subsection{Glide Ratio}
'Glide ratio' is a non-specific term which can refer to the ratio of horizontal to 
vertical motion with reference to either the surrounding airmass or the ground.

To reduce confusion, ground-referenced glide ratios (eg distance travelled over 
ground vs altitude lost) should be referred to by the term 'glide ratio over 
ground' when space allows, or 'glide ratio' / 'GR'.

Air-referenced glide ratios (eg airspeed vs sink rate) should be specified as 
'lift/drag ratio' / 'L/D ratio' / 'LD'. The lift/drag ratio is numerically equal 
to the air-referenced glide ratio when flying at constant speed.

If usage spans both air-referenced and ground-referenced glide ratios, the 
non-specific term 'glide ratio' / 'GR' should be used. 'Lift/drag ratio' should 
never be used to refer to ground-referenced glide ratios.

\chapter{Lua Scripting}\label{cha:lua}
\lstnewenvironment{lua}
{\lstset{language=[5.0]Lua}}
{}

Starting with version 7.0, \xc can be extended using
\href{http://www.lua.org/}{Lua} scripts.

Lua is a language that is easy to learn, powerful enough for XCSoar
and light: the interpreter library weighs just 200 kB.  Lua is a
common language choice for integrated scripting languages.

\section{Learning Lua}

\begin{lua}
print("Hello World")
\end{lua}

This manual will not attempt to teach you basic Lua.  There are enough
resources on the internet, for example:

\begin{itemize}
\item \href{http://www.lua.org/manual/5.3/}{Lua 5.3 Reference Manual}
\item \href{http://www.lua.org/pil/contents.html}{Programming in
  Lua}, a book on Lua
\item \href{http://lua-users.org/wiki/TutorialDirectory}{Tutorial
  Directory} on the lua-users wiki
\item \href{https://en.wikipedia.org/wiki/Lua_%28programming_language%29}{Wikipedia}
\end{itemize}

Just to get you started from here, here's some more example code:

\begin{lua}
-- comment starts with a double hyphen

--[[
multi
line
comment
]]--

i = 42
if i > 1 then
   print("i=" .. i)
elseif i == 0 then
   print("zero")
else
   error("negative")
end

a = {1, 'a', 3.14}
print(a[2])

function f(a, b)
   return a * b
end
print(f(2, 3))
\end{lua}

\section{Running Lua}

The directory \texttt{XCSoarData/lua/} may contain Lua scripts
(\texttt{*.lua}).  The directory \texttt{XCSoarData/lua/lib/} may
contain Lua libraries to be loaded with \texttt{require}.

After startup, \xc starts the script \texttt{init.lua} (if it
exists).

The \emph{InputEvent} ``\texttt{RunLuaFile}'' can be used to start
additional scripts.  If no parameter is given, the user is asked to
choose a file.  Note that the \emph{InputEvent} subsystem is
deprecated and will be removed once Lua support is complete.

As long as a Lua script runs, the \xc user interface is blocked.
Be careful not to write scripts that loop forever.

Once the Lua script finishes, the Lua interpreter is shut down --
unless the script has registered a callback (e.g. a \verb|timer|).  In
that case, the Lua script stays resident until it unregisters all
callbacks (or until \xc quits or the user stops the script
explicitly).

\section{Lua Standard Libraries}

\xc enables the following Lua standard libraries:

\begin{itemize}
\item \verb|package|
\item \verb|table|
\item \verb|string|
\item \verb|math|
\end{itemize}

Lua's \verb|print()| function writes to the \xc log file
(\texttt{XCSoarData/xcsoar.log}).

The \verb|error()| function aborts the Lua script and reports the
specified error message to the user.

\xc adds another function to the root namespace: \verb|alert()|.
It shows a dialog with the specified message, and returns as soon as
the user has closed the dialog.  This function is experimental, and
may disappear or be renamed at any time.  Most importantly: do not
abuse it, as it may annoy the user.

\section{\xc's Lua API}

The package/namespace \verb|xcsoar| provides access to \xc.  It
contains the following names:

\begin{maxipage}
\begin{tabularx}{1.9\textwidth}{l|X}
  Name & Description \\
  \hline
  \hline

  \verb|VERSION| & The \xc version number, for example
  ``\texttt{7.0}''. \\
  \hline
  \verb|blackboard| & Access to sensor data. \\
  \hline
  \verb|map| & The map view. \\
  \hline
  \verb|airspace| & Access to airspace data. \\
  \hline
  \verb|wind| & Access to wind data and settings. \\
  \hline
  \verb|logger| & Access to logger settings. \\
  \hline
  \verb|tracking| & Access to tracking settings. \\
  \hline
  \verb|timer| & Class for scheduling periodic callbacks. \\
\end{tabularx}
\end{maxipage}

\subsection{The Blackboard}

The blackboard provides access to sensor data, such as GPS location.

The following attributes are provided by \verb|xcsoar.blackboard|:

\begin{maxipage}
\begin{tabularx}{1.9\textwidth}{l|X}
Name & Description \\
\hline\hline

\verb|location| & The current location (table with
keys \verb|longitude| and \verb|latitude| in degrees) according to
GPS. \\

\hline

\verb|altitude| & The current altitude $[m]$ above MSL. \\

\hline

\verb|track| & The current flying direction above ground in degrees. \\

\hline

\verb|ground_speed| & The aircraft speed relative to the ground
$[{m \over s}]$. \\

\hline

\verb|air_speed| & The true airspeed
$[{m \over s}]$. \\

\hline

\verb|bank_angle| & The bank angle in degrees. \\

\hline

\verb|pitch_angle| & The pitch angle in degrees. \\

\hline

\verb|heading| & The current magnetic heading in degrees. \\

\hline

\verb|g_load| & The current g-load. \\

\hline

\verb|static_pressure| & The static pressure
$[{Pa}]$. \\

\hline

\verb|pitot_pressure| & The pitot pressure
$[{Pa}]$. \\

\hline

\verb|dynamic_pressure| & The dynamic pressure
$[{Pa}]$. \\

\hline

\verb|temperature| & The current temperature. \\

\hline

\verb|humidity| & The current humidity \\

\hline

\verb|voltage| & The external battery voltage
$[{V}]$. \\

\hline

\verb|battery_level| & The internal battery-level in percent. \\

\hline

\verb|noncomp_vario| & The non-compensated vertical speed
$[{m \over s}]$. \\

\hline

\verb|total_energy_vario| & The total-energy-compensated vertical speed
$[{m \over s}]$. \\

\hline

\verb|netto_vario| & The netto variometer value
$[{m \over s}]$. \\

\end{tabularx}
\end{maxipage}

Any of these may be \verb|nil| if its value is not known, e.g. if
there is no GPS fix.

\subsection{The Map}

The map provides access to XCSoar's map view.

The following attributes are provided by \verb|xcsoar.map|:

\begin{maxipage}
\begin{tabularx}{1.9\textwidth}{l|X}
Name & Description \\
\hline\hline

\verb|location| & The current reference location (may be aircraft location or
pan location). \\

\hline

\verb|is_panning| & Gives back if the panning mode is active at the moment.\\

\hline

\verb|enterpan()| & Activates the panning mode.\\

\hline

\verb|disablepan()| & Disables the panning mode.\\

\hline

\verb|leavepan()| & Leaves the panning mode.\\

\hline

\verb|panto(float latitude, float longitude)| & Pans to the given location.\\

\hline

\verb|pancursor(int dx, int dy)| & Pans the cursor by dx and dy. \\

\hline

\verb|zoom(int factor)| & Zooms the map, factor -2 to 2.\\

\hline

\verb|next()| & Opens the next Page.\\

\hline

\verb|prev()| & Opens the previous Page.\\

\hline

\verb|show()| & Show the map; disable thermal assistant or other
widgets replacing the map view. \\

\end{tabularx}
\end{maxipage}

\subsection{Airspace}

The Airspace provides access to airspace data, such as name / distance
to the next airspace.

The following attributes are provided by \verb|xcsoar.airspace|:

\begin{maxipage}
\begin{tabularx}{1.9\textwidth}{l|X}
Name & Description \\
\hline\hline

\verb|nearest_vertical_distance| & The vertical distance to the next airspace
$[{m}]$.\\

\hline

\verb|nearest_vertical_name| & The name of the next vertical airspace. \\

\hline

\verb|nearest_horizontal_distance| & The horizontal distance to the next airspace 
$[{m}]$. \\

\hline

\verb|nearest_horizontal_name| & The name of the next horizontal airspace. \\

\end{tabularx}
\end{maxipage}

\subsection{Task}

The Task provides access to task data such as distances / bearing to the 
next waypoint.

The following attributes are provided by \verb|xcsoar.task|:

\begin{maxipage}
\begin{tabularx}{1.9\textwidth}{l|X}
Name & Description \\
\hline\hline

\verb|bearing| & The true bearing to the next waypoint. 
For AAT tasks, this is the true \newline bearing 
to the target within the AAT sector. $[{degrees}]$\\

\hline

\verb|bearing_diff| & The difference between the glider's track bearing, 
to the bearing of \newline the next waypoint, or for AAT tasks, to the bearing 
to the target within \newline the AAT sector $[{degrees}]$.\\

\hline

\verb|radial| & The true bearing from the next waypoint 
to your position. $[{degrees}]$. \\

\hline

\verb|next_distance| & The distance to the currently selected waypoint. 
For AAT tasks, this \newline is the distance to the target within the AAT sector.
$[{m}]$ \\

\hline

\verb|next_distance_nominal| & The distance to the currently selected waypoint. 
For AAT tasks, this \newline is the distance to the origin of the AAT sector.
$[{m}]$ \\

\hline

\verb|next_ete| &  Estimated time required to reach next waypoint, 
assuming \newline performance of ideal MacCready cruise/climb cycle.\\

\hline

\verb|next_eta| & Estimated arrival local time at next waypoint, 
assuming performance \newline of ideal MacCready cruise/climb cycle. \\

\hline

\verb|next_altitude_diff| & Arrival altitude at the next waypoint relative 
to the safety arrival height. \\

\hline

\verb|nextmc0_altitude_diff| & Arrival altitude at the next waypoint with 
MC 0 setting, relative to the \newline safety arrival height. \\

\hline

\verb|next_altitude_require| & Additional altitude required to reach the next
turnpoint. \\

\hline

\verb|next_altitude_arrival| & Absolute arrival height at the next waypoint in 
final glide. \\

\hline

\verb|next_gr| & The required glide ratio over ground to reach the next waypoint, 
\newline given by the distance to the next waypoint divided by the height 
\newline required to arrive at the safety arrival height. \\

\hline

\verb|final_distance| & Distance to finish around remaining turn points. \\

\hline

\verb|final_ete| & Estimated time required to complete task, 
assuming performance \newline of ideal MacCready cruise/climb cycle. \\

\hline

\verb|final_eta| & Estimated arrival local time at task completion, 
assuming performance \newline of ideal MacCready cruise/climb cycle. \\

\hline

\verb|final_altitude_diff| & Arrival altitude at the final task turn point
relative to the safety arrival \newline height. \\

\hline

\verb|finalmc0_altitude_diff| & Arrival altitude at the final task turn point
, with MC 0 setting, relative \newline to the safety arrival height. \\

\hline

\verb|final_altitude_require| & Additional altitude required to finish 
the task. \\

\hline

\verb|task_speed| & Average cross country speed while on the current 
task, \newline not compensated for altitude. \\

\hline

\verb|task_speed_achieved| & Achieved cross country speed while on the 
current task, \newline compensated for altitude. Equivalent to 
Pirker cross country \newline speed remaining. \\

\hline

\verb|task_speed_instant| & Instantaneous cross country speed while 
on the current task, \newline  compensated for altitude. Equivalent 
to instantaneous Pirker cross \newline country speed. \\

\hline

\verb|task_speed_hour| & Average cross country speed while on the 
current task \newline over the last hour, not compensated for altitude. \\

\end{tabularx}
\end{maxipage}

\begin{maxipage}
\begin{tabularx}{1.9\textwidth}{l|X}
Name & Description \\
\hline\hline

\verb|final_gr| & The required glide ratio over the ground to finish 
the task, given by \newline the distance to go divided by the height 
required to arrive at the safety \newline arrival height. \\

\hline

\verb|aat_time| & Assigned Area Task time remaining. \\

\hline

\verb|aat_time_delta| & Difference between estimated task time and 
AAT miminum time.\\

\hline

\verb|aat_distance| & Assigned Area Task distance around target points 
for remainder of \newline task. \\

\hline

\verb|aat_distance_max| & Assigned Area Task maximum distance possible 
for remainder of \newline task. \\

\hline

\verb|aat_distance_min| & Assigned Area Task minimum distance possible 
for remainder of  \newline task\\

\hline

\verb|aat_speed| & Assigned Area Task average speed achievable around 
target points \newline remaining in minimum AAT time. \\

\hline

\verb|aat_speed_max| & Assigned Area Task average speed achievable 
if flying maximum \newline possible distance remaining in minimum 
AAT time. \\

\hline

\verb|aat_speed_min| & Assigned Area Task average spped achievable 
if flying minimum \newline possible distance remaining in minimum 
AAT time.  \\

\hline

\verb|time_under_max_height| & The contiguous period the plane has 
been below the task \newline start max. height. \\

\hline

\verb|next_etevmg| & Estimated time required to reach next waypoint, 
assuming current \newline ground speed is maintained. \\

\hline

\verb|final_etevmg| & Estimated time required to complete task, 
assuming current ground \newline speed is maintained. \\

\hline

\verb|cruise_efficiency| & Efficiency of cruse, 1 indicates perfect 
MacCready performance \\


\end{tabularx}
\end{maxipage}

\subsection{Settings}

The Settings provides access to xcsoar settings, such as MC value.

The following attributes are provided by \verb|xcsoar.settings|:

\begin{maxipage}
\begin{tabularx}{1.9\textwidth}{l|X}
Name & Description \\
\hline\hline

\verb|mc| & The current set MacCready Value
$[{m/s}]$.\\

\hline

\verb|bugs| & The current used bug settings in terms of polar degradation.\\

\hline

\verb|wingload| & The current wingload.\\

\hline

\verb|ballast| & Ballast of the glider. 0 means no ballst, 0.3 means 30\% of
the maximum \newline  ballast the glider can carry.\\

\hline

\verb|qnh| & Area pressure for barometric altimeter calibration
$[{Pa}]$.\\

\hline

\verb|max_temp| & The forecast ground temperature
$[{K}]$.\\

\hline

\verb|safetymc| & The MacCready setting used, when safety MC is enabled 
for reach \newline calculations, in task abort mode and for determining arrival altitude 
at \newline airfields.\\

\hline

\verb|riskfactor| & The STF risk factor reduces the MacCready setting used to 
calculate \newline speed to fly as the glider gets low, in order to  compensate for risk.\\

\hline

\verb|polardegradation| & A permanent polar degradation, 1 means no degradation,
0.5 indicates the \newline glider's sink rate is doubled.\\

\hline

\verb|arrivalheight| & The height above terrain that the glider should arrive
at for a safe landing.\\

\hline

\verb|terrainheight| & The height above trerrain that the glider must clear during
final glide.\\

\hline

\verb|setmc(float value)| &  Sets the MacCready value\\

\hline

\verb|setbugs(float value)| & Sets the bugs, 1.0 means no bugs, 0.5 means 50\% polar
degradation.\\

\hline

\verb|setqnh(float value)| & Sets the QNH $[{Pa}]$\\

\hline

\verb|setballast(float value)| & Sets the ballst, 0 means no ballst, 0.5 means 50\%
of the maximum \newline ballst the glider can carry.\\

\hline

\verb|setmaxtemp(float value)| & Sets the maximum temperature $[{K}]$.\\

\hline

\end{tabularx}
\end{maxipage}

\subsection{Wind}

The Settings provides access to xcsoar wind data and settings.

The following attributes are provided by \verb|xcsoar.wind|:

\begin{maxipage}
\begin{tabularx}{1.9\textwidth}{l|X}
Name & Description \\
\hline\hline

\verb|wind_mode| & Wind mode, 0: Manual, 1: Circling, 2: ZigZag, 3: Both.\\

\hline

\verb|setwindmode(int value)| & Sets wind mode $[0-3]$.\\

\hline

\verb|tail_drift| & Determines whether the snail trail is drifted with the 
wind \newline when displayed in circling mode, 0: Off, 1: On. \\

\hline

\verb|settaildrift(bool value)| & Turns Taildrift Off / On $[{0,1}]$.\\

\hline

\verb|wind_source| & The Source of the current wind, 0: None, 1: Manual,
\newline 2: Circling, 3: Both, 4: External.\\

\hline

\verb|wind_speed| & The current wind speed 
$[{m/s}]$.\\

\hline

\verb|setwindspeed(float value)| & Sets manual the wind speed $[{m/s}]$.\\

\hline

\verb|wind_bearing| & The current wind bearing 
$[{degrees}]$.\\

\hline

\verb|setwindbearing(float value)| & Sets manual the wind bearing $[{degrees}]$.\\

\hline

\verb|clear()| & Clears the wind settings and calculations.\\

\hline

\end{tabularx}
\end{maxipage}

\subsection{Logger}

The Settings provides access to xcsoar Logger data and settings.

The following attributes are provided by \verb|xcsoar.logger|:

\begin{maxipage}
\begin{tabularx}{1.9\textwidth}{l|X}
Name & Description \\
\hline\hline

\verb|pilot_name| & Gives back the set pilot name.\\

\hline

\verb|set_pilot_name(char* name)| &  Sets the pilot name.\\

\hline

\verb|time_step_cruise| & The time interval between logged 
points when not circling. $[s]$.\\

\hline

\verb|set_time_step_cruise(int time)| &  Sets time interval between logged 
points when not circling $[s]$.\\

\hline

\verb|time_step_circling| &  The time interval between logged 
points when circling $[s]$.\\

\hline

\verb|set_time_step_circling(int time)| &  Sets time interval between logged 
points when circling $[s]$.\\

\hline

\verb|auto_logger| &  Status of the auto-logger, 0: On, 1: Take off only
2: Off.\\

\hline

\verb|set_autologger(int mode)| &  Sets the Autologger mode, 0: On,
1: Take off only, 2: Off.\\

\hline

\verb|nmea_logger| &  Status of the NMEA-Logger, 0: Off, 1: On.\\

\hline

\verb|enable_nmea()| &  Enables the NMEA-Logger.\\

\hline

\verb|disable_nmea()| &  Disables the NMEA-Logger.\\

\hline

\verb|log_book| &  Status of the log-book, 0: Off, 1: On.\\

\hline

\verb|enable_logbook()| &  Enables the logbook.\\

\hline

\verb|disable_logbook()| &  Disables the logbook.\\

\hline

\verb|logger_id| &  The current set logger-id.\\

\hline

\verb|set_logger_id(char* id)| &  Sets the logger-id.\\

\end{tabularx}
\end{maxipage}

\subsection{Tracking}

The Settings provides access to xcsoar Tracking settings.

The following attributes are provided by \verb|xcsoar.tracking|:

\begin{maxipage}
\begin{tabularx}{1.9\textwidth}{l|X}
Name & Description \\
\hline\hline

\verb|skylines_enabled| & States if skylines tracking is enabled.\\

\hline

\verb|enable_skylines()| & Enables skylines tracking.\\

\hline

\verb|disable_skylines()| & Disables skylines tracking.\\

\hline

\verb|skylines_roaming| & States if skylines roaming is enabled.\\

\hline

\verb|skylines_interval| & The skylines tracking interval $[s]$.\\

\hline

\verb|set_skylines_interval(int interval)| & Sets the tracking interval $[s]$.\\

\hline

\verb|skylines_traffic_enabled| & If enabled shows friends on the map,
download the position \newline of your friends live from the SkyLines server.\\

\hline

\verb|enable_skylines_traffic()| & Enables the display of friends on the map.\\

\hline

\verb|disable_skylines_traffic()| & Disables the display of friends on the map.\\

\hline

\verb|skylines_near_traffic_enabled| & If enabled shows nearby traffic.\\

\hline

\verb|enable_skylines_near_traffic()| & Enables the display of nearby traffic on the map.\\

\hline

\verb|disable_skylines_near_traffic()| & Disables the display of nearby traffic on the map.\\

\hline

\verb|livetrack24_enabled| & States if livetrack24 is enabled.\\

\hline

\verb|enable_livetrack24()| & Enables livetrack24.\\

\hline

\verb|disable_livetrack24()| & Disables livetrack24.\\

\hline

\verb|livetrack24_interval| & Livetrack24 tracking interval $[s]$.\\

\hline

\verb|set_livetrack24_interval(int interval)| & Sets the tracking interval $[s]$.\\

\hline

\verb|livetrack24_vehicle_name| & Set current set vehicle name.\\

\end{tabularx}
\end{maxipage}

\subsection{Timers}

The class \verb|xcsoar.timer| implements a timer that calls a given
Lua function periodically.

\begin{lua}
xcsoar.timer.new(60, function(t)
  print("A minute has passed")
  t:cancel()
end)
\end{lua}

The following methods are available in \verb|xcsoar.timer|:

\begin{maxipage}
\begin{tabularx}{1.9\textwidth}{l|X}
Method & Description \\
\hline\hline

\verb|new(period, function)| & Create a new instance and schedule
it.  The period is a numeric value in seconds. \\

\hline

\verb|cancel()| & Cancel the timer. \\

\hline

\verb|schedule(period)| & Reschedule the timer. \\

\end{tabularx}
\end{maxipage}

\subsection{Legacy}

Before version 7.0, \xc was adapted using the \emph{InputEvent}
subsystem (see Appendix~\ref{cha:file_formats}).  During the Lua
transition, Lua scripts can invoke InputEvents, for example:

\begin{lua}
xcsoar.fire_legacy_event("Setup", "basic")
xcsoar.fire_legacy_event("Zoom", "basic")
\end{lua}

This function will be removed before the final 7.0 release.


\chapter{File formats}\label{cha:file_formats}

\section{Input Events}

\subsection{Introduction}

The Input System is actually a large number of things all bunched into one.

Primarily it is about giving the user control of what button does what
and when. There is a new concept called Input Mode - this is a the
mode the GUI is in for input. For example, you can click on the info
boxes and you are now in "infobox" mode. Clicking on the map is called
"default". But it doesn't stop there, you can create a new mode called
anything you like. This may not mean much - but wait till you combine
it with the rest of the features...

Input is not restricted to hardware buttons any more. You can map all
your hardware buttons (currently support for APP1 to APP6, Left,
Right, Up, Down and Enter, although I believe we can do some more) but
also any key code at all. This feature allows those with a built in
keyboard to use any key to map to any function in XCS. Where it comes
into real advantage is in external keyboards. There are a number of
bluetooth devices out there (eg: \url{http://shop.brando.com.hk/btgamepad.php}) 
which can map each of their
buttons to any key code - that key code can then be mapped to any
feature in XCS. You can then add to the hardware buttons the buttons
available to you on external devices. Other inputs (eg: Serial) are
also being looked at - and support is in the code for that extension.

We are striving towards a platform which is not only easier to use and
more intuitive, but also faster and easier to use in flight as
well. As such, another new feature as part of input is the concept of
Button Labels. Combined with the modes mentioned above, you can create
any arbitrary set of functions to map to any number of buttons. Think
about it like creating a tree, or a multiple level menu.

This produces two benefits that I know will be appreciated by people
with limited inputs. The first is that you can create menus, where by
you press one button to get to the next level (eg: pressing on APP1
brings up AutoZoom, Pan Mode, Full screen on the other buttons. Press
APP1 again and it goes to Terrain, Marker and Auto MacCready. Press
APP1 again and the menu is gone) - but more importantly for those with
touch screens and limited buttons, each of these labels can optionally
be assigned a key and you can touch the button area as if it was a
button.  This means that we can actually control on a touch screen
model the entire system without buttons - press an area of the screen
and the buttons pop up, click through - change options and more.

The combined features of labels, configurable buttons (including from
external hardware), hierarchical menus (for lack of a better name),
touch screen buttons has allowed us to configure XCS - without
recompile - for an enormous range of hardware, and personal
preference.  And all configurable as plane text, simple files. There
is no need for a file, the defaults internally will probably be a
combination of a 4 button bottom system with one button always shown
on screen for no/few button display.

The screen layout - location of the labels - is also totally
configurable - allowing us to vary the layout of buttons depending on
the type of organiser or desired look and feel.

There is a great unexpected benefit in the development of the input
system.

We can execute any number of events attached to an input with only 2
extra lines of code. This worked perfectly. So now we have a basic
macro system, allowing many more events to be attached to a single
input event.

But it doesn't stop there, this has lead to some more excellent
developments. The idea of Glide Computer Events things like "Maximum
Height Reached". Currently we play a sound effect for that. But you
may choose to play a sound, bring up a message box and write to the
log file.

One nice feature of XCS is the ability to change things such as Zoom
and North when Circling. Now you can do so much more. You could choose
to point North, Zoom to 1.0 (rather than a relative change), Turn on
Vario Sounds, Start a timer. When switching back to Cruise mode, you
can bring up the stats box for 30 seconds. The options are limited by
your imagination.

This is also contributing to a major reduction in complex code. We can
move out these complex tests into one centrally, easier to manage
system, reducing bugs and improving maintainability.

Another side benefits of these Macros is User Defined Flight
Modes. One idea was a button which switched to Zoom 1.0, Pan ON, Pan
Move to Next Waypoint. Basically the ability to jump and see the next
waypoint. And in the previous we can change the Input Mode to
"ViewWaypoint" - at which point you can redefine the same button to
switch back to your original settings.

The flexibility of this system comes with only one small price. We
can't provide an interface within XCS to fully customise all of these
near infinitely variable possibilities. However I believe that is
unnecessary anyway, you are not likely to change these sort of
features very often, and definitely not on the field. That does not
mean you can't, you can of course edit the plane(sic) text file to
change functions.

What this really means is that we can have people in the project
helping and contributing to the customising of XCS, without having to
change the code. This, especially on an open source project is
fantastic as it nicely separates the user interface changes from the
highly reliable part of the code. It also involves people who can
develop new interfaces and functions that are expert gliders but not
necessarily programmers.

For information on file formats see Common/Data/Input/template.xci and
the web site documentation.


\subsection{Defaults and Files}

The file in the source Common/Data/input/template.xci is used to generate
automatically the C code necessary for the default configuration. However it is
in the exact same format as can be read in by XCS and therefore can be used
literally as a template for a more complicated file.

When you create your own file, you will need to select it as the Input File
within XCSoar Menu/Settings/Input Files, and then restart XCS.


\subsection{File format}

The file is plain text, with key=value pairs and a blank line to indicate
the end of a record.

\begin{verbatim}
mode=default
type=key
data=APP1
event=StatusMessage My favorite settings are done
event=ScreenModes full
event=Sounds on
event=Zoom 1.0
event=Pan off
label=My Prefs
location=1
\end{verbatim}

The record above demonstrates remapping the first hardware key on your
organiser to change Pan to off, Zoom to 1.0 Sounds on, ScreenModes
full, and then a status message to tell you it is done.

Lines are terminated by the stanard DOS newline which is CRLF (Carrage
Return then Line Feed). Records are terminated by an extra new line.

\subsection{Event order}

Until further work is done on processing, events are actually done in
reverse order - also known as RPN. This is because the events work on
the stack principle. Each one is pushed onto the stack for execution,
and then executed by popping back off the stack. This has reduced
complexity of the code base.

When writing input events, have a look where you put the StatusMessage
and make sure that it is at the top, not the bottom (if you have one).

\subsection{Event list}

{\footnotesize
\begin{tabular}{l|p{7cm}}
\emph{Event} & \emph{Description} \\

\hline

\texttt{ChangeInfoBoxType} \emph{C} & Possible arguments: next,
previous. \\

\hline

\texttt{MainMenu} & \\

\hline

\texttt{MarkLocation} & Mark a location. \\

\hline

\texttt{Mode} \emph{M} & Set the screen mode. \\

\hline

\texttt{Pan} \emph{P} & Control pan mode.  Possible arguments: on
(enable pan), off (disable pan), up, down, left, right \\

\hline

\texttt{PlaySound} \emph{S} & Play the specified sound. \\

\hline

\texttt{SelectInfoBox} \emph{S} & Possible
arguments: next, previous. \\

\hline

\texttt{SnailTrail} \emph{S} & Change snail trail setting.  Possible
arguments: off, short, long, show. \\

\hline

\texttt{ScreenModes} \emph{M} & Set the screen mode.  Possible
arguments: normal, auxilary, toggleauxiliary, full, togglefull,
toggle. \\

\hline

\texttt{Sounds} \emph{S} & Change vario sounds.  Possible arguments:
toggle, on, off, show. \\

\hline

\texttt{StatusMessage} \emph{M} & Display the specified status
message. \\

\hline

\texttt{Zoom} \emph{Z} & Everything about zoom of map.  Possible
arguments: auto toogle, auto on, auto off, auto show, in, out, +, ++,
-, --. \\

\end{tabular}}

\subsection{Modes}

XCSoar now has the concept of Modes. These are an arbitrary
string that associates with where and what XCS is doing.

Note: a mode entry in a record can have multiple entries by using a space
between eg: "infobox menu1 menu2"

\subsubsection{List of known modes}

\begin{description}
\item[\texttt{default}] Really map mode, where you mostly are
\item[\texttt{infobox}] An info box has been selected on the scrreen
\item[\texttt{*}] Any other arbitrary string
\end{description}

Mode precedence has been tricky, so instead of solving the problem 
it is being worked around. XCS will choose to set a global variable 
to specify what mode it thinks it is in. This can then be used by the
input code to decide what to do. This mode could get out of sink
with the real world, and careful checking will be required, but at
this stage it seems like the only sensible option.

The code will review first if an entry exists in the current mode, and 
then in the default mode. This allows you to do one of the following
example: Define a default action for button "A" to be "Zoom In" but
make that button increase Bugs value in infobox mode only. You can do
this by making an "default" and a "infobox" entry. You can also put an entry
in for Button "A" for every mode and have complete control.

Special Modes - eg: the level of a menu (Think File vs Edit, vs Tools vs Help)

have special modes, such as
the level of the menu you are at. You press one button, then another
set become available (like pressing menu and seeing Settings etc). This
will be very useful in non-touch screen models. The menu configuration
can then be read from this same file and configured, allowing any
number of levels and any number of combinations.

The only hard part is what mode to go back to. We need a 
"Calculate Live Mode" function - which can be called to calculate the
real live mode (eg: finalglide vs curse) rather than the temporary
mode such as Menu, Special Menu Level, Warning etc.

The label and location values are examples of what can be done here
to allow input button labels to be displayed. What needs to be 
considered is a simple way of mapping the locations and the size.
In some models it may be that buttons are 4 across the top of the
screen, where as others it is 3 or 2 or even 6. So both size and
location needs to be considered.

The label itself will go through gettext to allow language
translations.

\subsection{Keys}

The key type can have the following possible values:

\begin{description}
\item[\texttt{APP1-APP6}] Hardware key on pocket pc
\item[\texttt{F1-F12}] Standard function keys
\item[\texttt{LEFT, RIGHT, UP, DOWN, RETURN}] Mapped to arrow keys - joystick on
  organisers
\item[\texttt{A-Z, 0-9}] and other possible keyboard buttons (case is ignored)
\end{description}


XXX Review...
Input Types

Types:

hardware	These are the standard hardware buttons 
on normal organisers. Usually these are
APP1..6.

keyboard	Normal characters on the keyboard (a-z etc)

nmea		A sentence received via NMEA stream (either)

virtual		Virtual buttons are a new idea, allowing
multiple buttons to be created on screen.
These buttons can then be optionally mapped
to physical buttons or to a spot on the 
screen (probably transparent buttons over 
the map).

Modifiers

It is a long term goal of this project to allow modifiers for keys.
This could include one of the following possibilities:

\begin{itemize}
\item Combination presses (although not supported on many devices)
\item Double Click
\item Long Click
\end{itemize}

Modifiers such as the above will not be supported in the first release.

\begin{verb}
Functions/Events - what it does

AutoZoom		on, off, toggle 
FullScreen		on, off, toggle
SnailTrail 		on, off, long, toggle
VarioSound 		on, off
Marker 			optional text to add
MenuButton 		on, off, toggle
Menu			open, close, toggle
MenuEntry		task, b+b, abortresume, abore, resume, pressure
logger, settings, status, analysis, exit, cancel
NOTE: Some of the above may be separate functions
Settings		(each setting, bring up to that point)
Bugs			add, subtract, 0-100% (set value)
Ballast			add, subtract, 0-100% (set value)
Zoom			add, subtract, 0-nn (set value)
Wind			up, down, 0-nn (set value, left, right, "n","ne","e","se","s","sw","w","nw"...
MacCready		add, subtract, 0-nn (set value)
WaypointNext		"String" to specific waypoint
eg: WayPointNext "home"
WayPoint???		"reverse" - reverse, from last passed back to start (ie: from here to home)
"drop next" - drop the next
"restore" - restore all - from start of flight but 
XXX This needs more thought
flight 			"startstop", "start", "stop", "release"
Start/Stop of flight - Can be automatic, but pressing will override
automatic part.
release 		markse the point of release from tow
\end{verb}

\subsection{Glide Computer Events}

These are automatically triggered events. They work in exactly the
same way, but instead of the user pressing a key, the glide computer
triggers the events.

A simple example is moving from Cruise to Climb mode. We want to zoom
in, change our track up to north up and switch to full screen. You may
also choose to drop a marker with the words "entered thermal". The
choicese are up to your imaginations - the GCE (Glide Computer Events)
allow you to control what happens.

These are represented as "type=gce" and data=* - as listed below.

{\footnotesize
\begin{tabular}{l|p{5.6cm}}
\emph{Event} & \emph{Description} \\
\hline

\verb|COMMPORT_RESTART| & The comm port is restarted. \\

\hline

\verb|FLIGHTMODE_CLIMB| & The flight mode has switched to ``climb''. \\

\hline

\verb|FLIGHTMODE_CRUIS| & The flight mode has switched to ``cruise''. \\

\hline

\verb|FLIGHTMODE_FINALGLIDE| & The flight mode has switched to ``final
glide''. \\

\hline

\verb|GPS_CONNECTION_WAIT| & Waiting for the GPS connection. \\

\hline

\verb|GPS_FIX_WAIT| & Waiting for a valid GPS fix. \\

\hline

\verb|HEIGHT_MAX| & Maximum height reached for this trip. \\

\hline

\verb|LANDING| & You are at landing. \\

\hline

\verb|STARTUP_REAL| & First message - this happens at startup of the
real XCS. \\

\hline

\verb|STARTUP_SIMULATOR| & Startup first message. This happens during
simulator mode. \\

\hline

\verb|TAKEOFF| & You have taken off. \\

\end{tabular}}

\section{Map Data file formats}
The map data is typically downloaded from the map generator server and consists of a single \texttt{.xcm} file.
It is a zip file which contains several separate files for terrain, topography and waypoint data:
\begin{itemize}
\item[\texttt{info.txt}] General map information
\item[\texttt{terrain.jp2}] Terrain (elevation) data, georeferenced in \texttt{terrain.j2w}
\item[\texttt{waypoints.cup}] Waypoint data
\item[\texttt{topology.tpl}] Topography data file index
\item[\texttt{*.shp} / \texttt{*.dbf} / \texttt{*.shx}] A set of ESRI format shape file sets with actual topography voctor data information, as listed and defined in \texttt{topology.tpl} (Coasts, rivers, roads, cities etc.)
\end{itemize}
\subsection{Map information}
\texttt{info.txt}
Contains information about the map as a whole, such as creator, creation time, and lat/lon range.

\subsection{Terrain data files}
The map cointains a digital elevation model of the map area. It is stored as an JPEG2000 compressed image in the file \texttt{terrain.jp2}. The projection information (lat/lon boundaries) of the DEM file are contained in the text file \texttt{terrain.j2w}, in decimal degree latitude/longitude format.
Water is defined as elevation lower than \texttt{TERRAIN\_WATER\_THRESHOLD=-30000}, therefore care has to be taken that JPEG compression parameters and algorithms  are used which do not generate artefacts at the coastlines due to the potentially big jump in elevation value.
\subsection{Waypoints}
A map database file can contain waypoints. They reside in the \texttt{waypoints.cup} file, which has regular \texttt{.CUP} format.

\subsection{Topography data}
\subsubsection{Shape files}
Non-elevation topography data is stored in standard ESRI shape files. For each type of topographic shape (road, river, city outline, etc.) there is one shape filein \texttt{.shp}, which containes all shapes of this type. For each \texttt{.shp} file, there has to be an associated \texttt{.dbf} file containing shape metadata (such as the name of the city) in dBASE format, and an index file of \texttt{.shx} file type which contains the index that relates the metadata to the shapes. 

All of this is defined in the ESRI shape file standard. The official definition of the standard can be found at\\ \href{http://www.esri.com/library/whitepapers/pdfs/shapefile.pdf}{http://www.esri.com/library/whitepapers/pdfs/shapefile.pdf}, but there are more compact descriptions available on the web, see for example wikipedia info and links at \\ \href{http://en.wikipedia.org/wiki/Shapefile}{http://en.wikipedia.org/wiki/Shapefile}.

There can be more files associated with each shape file, such as \texttt{.prj}, \texttt{.qix}, \texttt{.atx}, which are not used by XCSoar.

The set of shapefiles actually used by XCSoar and the attributes of each file are defined in the topography layer description file \texttt{topology.tpl}. All shape files used by the map must be listed there.

\subsubsection{Topography layer description file (topology.tpl) format}

Each line of the topography layer description file (\texttt{topology.tpl}) contains
a comma separated list (CSV) of information needed for rendering of an
individual topography layer. Lines starting with '*' are ignored.

XCSoar v6.6 and earlier will display at most 20 topography layers.
XCSoar v6.7 and later will display at most 30 topography layers.
\begin{table}[ht]
\centering
\sffamily
\begin{tabular}{@{}lll@{}}
\toprule
\addlinespace
Column name&Data type&Valid range\\
%\midrule
\cmidrule(r){1-1}\cmidrule(lr){2-2}\cmidrule(l){3-3}
filename&string&\\
range&double (nm)&-\\
icon&string&\\
label index&int&0-1\\
color (red component)&int&0-255\\
color (green component)&int&0-255\\
color (blue component)&int&0-255\\
pen width&int&0-31\\
label range (nm)&double&-\\
important label range (nm)&double&-\\
alpha&int&0-255\\
%\addlinespace
\bottomrule
\end{tabular}
\caption{Topography file format}
\label{tab:topography-file-format}
\end{table}

\begin{description}
\item[filename] The filename of the Topography layer within the container file.
\item[icon] XCSoar v6.5 and earlier, Only the value 219 is recognised, for town icons.
From XCSoar v6.6, the name of the icon to display. Optional. See below for a list of available names.
\item[range] Zoom level threshold. Layer elements will not be drawn unless zoomed in closer than this threshold.
\item[pen width] Lines contained within this layer are drawn with pen width.
\item[label range] Label display zoom level threshold. Labels contained in the layer file
will not be rendered unless zoomed in closer than this threshold.
\item[important label range] A zoom level threshold. Labels contained in the layer file will be
rendered in standard style when the display zoom level is greater than this threshold.
\item[alpha] The alpha component controls transparency of polygons...
0 means polygons are completely transparent, 255 means they are completely opaque.
Only used by XCSoar v6.7 and later.

\texttt{Versions of XCSoar running on Windows and WinCE ignore any item where transparency is specified.}

\end{description}
\subsubsection{Point Features}
Prior to XCSoar v6.6, this could contain the value 219 to display an icon for a town
From XCSoar v6.6, a user can put an optional string into the icon column in topology.tpl in the .XCM file (e.g.)
\begin{itemize}
\item SpotHeight,5,\texttt{mountain\_top},1,64,64,64,1,5,
\item Mast,10,\texttt{obstacle},,,,,1,10,
\end{itemize}
This can be used for Shapefiles containing point features or polygons or linestrings, but is probably only useful for point features.

The icon of the corresponding image and optional label will be displayed. In the first example, 
the ``mountain\_top'' icon and a label will be displayed for each point in the SpotHeight shapefile. My
SpotHeight Shapefile has been generated with the point elevation in feet as the label value). For the second example, only ``obstacle'' icons
(no labels) will be displayed for points in the Mast Shapefile..

Icon names are detected in TopographyStore.cpp. Names must be given in lowercase. If the icon name given is unknown, or no icon name is given, then icons are not displayed for that Shapefile.


Names correspond to images which have been linked into XCSoar, although
it is envisaged that in future these will be names of icon files. Available icon names are:
\begin{itemize}
\item mountain\_top \includegraphics[angle=0,width=8pt,keepaspectratio='true']{icons/map_mountain_top.pdf}
\item bridge \includegraphics[angle=0,width=8pt,keepaspectratio='true']{icons/map_bridge.pdf}
\item tunnel \includegraphics[angle=0,width=8pt,keepaspectratio='true']{icons/map_tunnel.pdf}
\item tower \includegraphics[angle=0,width=8pt,keepaspectratio='true']{icons/map_tower.pdf}
\item power\_plant \includegraphics[angle=0,width=8pt,keepaspectratio='true']{icons/map_power_plant.pdf}
\item obstacle \includegraphics[angle=0,width=8pt,keepaspectratio='true']{icons/map_obstacle.pdf}
\item mountain\_pass \includegraphics[angle=0,width=8pt,keepaspectratio='true']{icons/map_pass.pdf}
\item weather\_station \includegraphics[angle=0,width=8pt,keepaspectratio='true']{icons/map_weather_station.pdf}
\item thermal\_hotspot
\item town
\item mark \includegraphics[angle=0,width=8pt,keepaspectratio='true']{icons/map_flag.pdf}
\item turnpoint \includegraphics[angle=0,width=8pt,keepaspectratio='true']{icons/map_turnpoint.pdf}
\item small
\item cruise \includegraphics[angle=0,width=8pt,keepaspectratio='true']{icons/mode_cruise.pdf}
\item terrainwarning
\item logger
\item loggeroff
\item target
\item teammate\_pos
\item airspacei
\item traffic\_safe
\item traffic\_warning
\item traffic\_alarm
\item taskturnpoint
\item marginal \includegraphics[angle=0,width=8pt,keepaspectratio='true']{icons/winpilot_marginal.pdf}
\item landable \includegraphics[angle=0,width=8pt,keepaspectratio='true']{icons/winpilot_landable.pdf}
\item reachable \includegraphics[angle=0,width=8pt,keepaspectratio='true']{icons/winpilot_reachable.pdf}
\item airport\_reachable \includegraphics[angle=0,width=8pt,keepaspectratio='true']{icons/alt_reachable_airport.pdf}
\item airport\_unreachable \includegraphics[angle=0,width=8pt,keepaspectratio='true']{icons/alt_landable_airport.pdf}
\item airport\_marginal \includegraphics[angle=0,width=8pt,keepaspectratio='true']{icons/alt_marginal_airport.pdf}
\item airport\_unreachable2 \includegraphics[angle=0,width=8pt,keepaspectratio='true']{icons/alt2_landable_airport.pdf}
\item airport\_marginal2 \includegraphics[angle=0,width=8pt,keepaspectratio='true']{icons/alt2_marginal_airport.pdf}
\item outfield\_unreachable2 \includegraphics[angle=0,width=8pt,keepaspectratio='true']{icons/alt2_landable_field.pdf}
\item outfield\_marginal2 \includegraphics[angle=0,width=8pt,keepaspectratio='true']{icons/alt2_marginal_field.pdf}
\item outfield\_reachable \includegraphics[angle=0,width=8pt,keepaspectratio='true']{icons/alt_reachable_field.pdf}
\item outfield\_unreachable \includegraphics[angle=0,width=8pt,keepaspectratio='true']{icons/alt_landable_field.pdf}
\item outfield\_marginal \includegraphics[angle=0,width=8pt,keepaspectratio='true']{icons/alt_marginal_field.pdf}
\end{itemize}
\subsubsection{Adding new Icons}
At the moment, adding new icons requires a rebuild of the XCSoar application.It is envisaged that, in future, 
this process won't be required\dots users will include icon files in their .XCM map container files, and refer to them
by name. However, that has not yet been implemented. 

To add your own images to the list of icons:
\begin{enumerate}
\item Create a .svg file for the icon (e.g.\ mast.svg) and copy into \texttt{xcsoar/Data/icons}. For Android, the name must be lowercase.
\item Insert two (for normal and high-res) lines into \texttt{xcsoar/Data/XCSoar.rc},  (e.g.)
\begin{verbatim}
BITMAP\_ICON(IDB\_MAST, "mast")
BITMAP\_ICON(IDB\_MAST\_HD, "mast\_160")
\end{verbatim}
\item Insert two lines into \texttt{xcsoar/src/Resources.hpp} (e.g.)
\begin{verbatim}
MAKE_RESOURCE(IDB_MAST, 500);
MAKE_RESOURCE(IDB_MAST_HD, 5500);
\end{verbatim}
\item Add a corresponding line into the \texttt{icon\_list} table in \texttt{xcsoar/src/Topography/TopographyStore.cpp}
\begin{verbatim}
  {"mast", IDB\_MAST},
\end{verbatim}
\item Make XCSoar
\end{enumerate} 
After this, a line can be added in \texttt{topology.tpl} to connect the icon to the Shapefile using the icon name. (e.g.)
\begin{verbatim}
Mast,10,mast,,,,,1,10,
\end{verbatim}

Note that unless these changes are merged into the main XCSoar repository, then only your specific build of XCSoar will be able to 
display your icon image.



%%%%%%%%%%%%%%%%%%%%%%%%%%%%%%%%%%%%%%%%%%%%%%%%%%%
\appendix

\chapter{Setting up a development environment based on linux}\label{cha:developmentsetup}
This describes the setup of a development environment suitable to compile \xc for most supported platforms. The manual focuses on debian-based flavors of linux, but should be easily adapted to other distributions. It is developed with Ubuntu 16.04 LTS, and was last changed in May 2016 for \xc 7.0 preview  6.

\section{Install base system}\label{sec:developsetup-base}
The procedure uses the standard install, during/after which you should update all packages to the current level using

\begin{verbatim}
sudo apt-get update
sudo apt-get upgrade
\end{verbatim}

\section{Platform-independent components}\label{sec:developsetup-general}

Components required for all platforms, plus git for downloading the source and uploading commits.

\begin{verbatim}
sudo apt-get install make \
  librsvg2-bin xsltproc \
  imagemagick gettext ffmpeg \
  git quilt
\end{verbatim}
Some of these packages might be already present in your distribution (e.g. in the case of Ubuntu).

It is also recommended (though not required) to install \texttt{ccache}, which can speed up building a lot, when rebuilding after a \texttt{clean}.
\begin{verbatim}
sudo apt-get install ccache
\end{verbatim}


\section{Download source code}\label{sec:developsetup-sourcecode}

Download the source code of \xc by executing /textt{git} in the following way in your project directory
\begin{verbatim}
git clone git://git.xcsoar.org/xcsoar/master/xcsoar.git

# download submodule
cd xcsoar
git submodule init
git submodule update
\end{verbatim}

\section{For Unix builds}

Make sure your distribution contains a gcc compiler version 4.9 or higher in order to build the Unix target.

In addition, the following packages are needed:

\begin{verbatim}
sudo apt-get install make g++ \
  zlib1g-dev \
  libsdl1.2-dev libfreetype6-dev \
  libpng-dev libjpeg-dev \
  libtiff5-dev libgeotiff-dev \
  libcurl4-openssl-dev \
  liblua5.2-dev lua5.2-dev \
  libxml-parser-perl \
  librsvg2-bin xsltproc \
  imagemagick gettext \
  libegl1-mesa-dev \
  fonts-dejavu
\end{verbatim}

\section{For Android  builds}
\subsection{Install JDK}
First of all, we need Java to work with Android. We use OpenJDK 8 in this example, but OpenJDK 7 should also work.

\begin{verbatim}
sudo apt-get install openjdk-8-jdk ant
sudo apt-get install vorbis-tools automake
\end{verbatim}

\subsection{Install Android SDK and NDK}

The SDK and NDK are installed in the user directory. The \texttt{Makefile} expects them in specific directories, therefore we create some symlinks.

\begin{verbatim}
cd ~/
mkdir opt
cd opt/

# Download and extract SDK
wget "http://dl.google.com/android/android-sdk_r23-linux.tgz"

tar xvfz android-sdk_r23-linux.tgz
mv android-sdk-linux/ android-sdk-linux_x86/
ln -s android-sdk-linux/ android-sdk-linux_x86
rm android-sdk_r23-linux.tgz

# Download NDK installer

wget "http://dl.google.com/android/\
repository/android-ndk-r12b-linux-x86_64.zip"

unzip android-ndk-r12b-linux-x86_64.zip
\end{verbatim}

We need to add the paths by appending the following lines to \texttt{\textasciitilde/.profile} and \texttt{\textasciitilde/.bashrc} --- here relative paths work.

\begin{verbatim}
# Android SDK
export ANDROID_SDK_HOME="~/opt/android-sdk-linux_x86"
export PATH="$PATH:$ANDROID_SDK_HOME/tools"
export PATH="$PATH:$ANDROID_SDK_HOME/platform-tools"

# Android NDK
export ANDROID_NDK_HOME="~/opt/android-ndk-r11"
export PATH="$PATH:$ANDROID_NDK_HOME"
\end{verbatim}

The Android SDK (which contains some 32bit components) seems to have trouble running on some 64bit distributions, the following packages fix this:

\begin{verbatim}
sudo apt-get install lib32z1 lib32ncurses5 \
libgcc1:i386
\end{verbatim}

Now, you should start the Android manager with the \texttt{android} command, and install the required packages. Use the GUI to select and install SDK 22, and the (preselected) newest tools. Be warned, these will need many GB of space on your harddisk! (approximately 15 GB for typical defaults)

Finally, for the \texttt{adb} debugger interface to work via USB, most distributions require manual fixing of the \texttt{udev} rules. This can be done in the following way:

\begin{maxipage}
\begin{verbatim}
#Create etc/udev/rules.d/99-android.rules file

echo "#Acer" > 99-android.rules
echo "SUBSYSTEM=="usb", SYSFS{idVendor}=="0502", MODE="0666"" >> 99-android.rules
echo "#ASUS" >> 99-android.rules
echo "SUBSYSTEM=="usb", SYSFS{idVendor}=="0b05", MODE="0666"" >> 99-android.rules
echo "#Dell" >> 99-android.rules
echo "SUBSYSTEM=="usb", SYSFS{idVendor}=="413c", MODE="0666"" >> 99-android.rules
echo "#Foxconn" >> 99-android.rules
echo "SUBSYSTEM=="usb", SYSFS{idVendor}=="0489", MODE="0666"" >> 99-android.rules
echo "#Garmin-Asus" >> 99-android.rules
echo "SUBSYSTEM=="usb", SYSFS{idVendor}=="091E", MODE="0666"" >> 99-android.rules
echo "#Google" >> 99-android.rules
echo "SUBSYSTEM=="usb", SYSFS{idVendor}=="18d1", MODE="0666"" >> 99-android.rules
echo "#HTC" >> 99-android.rules
echo "SUBSYSTEM=="usb", SYSFS{idVendor}=="0bb4", MODE="0666"" >> 99-android.rules
echo "#Huawei" >> 99-android.rules
echo "SUBSYSTEM=="usb", SYSFS{idVendor}=="12d1", MODE="0666"" >> 99-android.rules
echo "#K-Touch" >> 99-android.rules
echo "SUBSYSTEM=="usb", SYSFS{idVendor}=="24e3", MODE="0666"" >> 99-android.rules
echo "#KT Tech" >> 99-android.rules
echo "SUBSYSTEM=="usb", SYSFS{idVendor}=="2116", MODE="0666"" >> 99-android.rules
echo "#Kyocera" >> 99-android.rules
echo "SUBSYSTEM=="usb", SYSFS{idVendor}=="0482", MODE="0666"" >> 99-android.rules
echo "#Lenevo" >> 99-android.rules
echo "SUBSYSTEM=="usb", SYSFS{idVendor}=="17EF", MODE="0666"" >> 99-android.rules
echo "#LG" >> 99-android.rules
echo "SUBSYSTEM=="usb", SYSFS{idVendor}=="1004", MODE="0666"" >> 99-android.rules
echo "#Motorola" >> 99-android.rules
echo "SUBSYSTEM=="usb", SYSFS{idVendor}=="22b8", MODE="0666"" >> 99-android.rules
echo "#NEC" >> 99-android.rules
echo "SUBSYSTEM=="usb", SYSFS{idVendor}=="0409", MODE="0666"" >> 99-android.rules
echo "#Nook" >> 99-android.rules
echo "SUBSYSTEM=="usb", SYSFS{idVendor}=="2080", MODE="0666"" >> 99-android.rules
echo "#Nvidia" >> 99-android.rules
echo "SUBSYSTEM=="usb", SYSFS{idVendor}=="0955", MODE="0666"" >> 99-android.rules
echo "#OTGV" >> 99-android.rules
echo "SUBSYSTEM=="usb", SYSFS{idVendor}=="2257", MODE="0666"" >> 99-android.rules
echo "#Pantech" >> 99-android.rules
echo "SUBSYSTEM=="usb", SYSFS{idVendor}=="10A9", MODE="0666"" >> 99-android.rules
echo "#Philips" >> 99-android.rules
echo "SUBSYSTEM=="usb", SYSFS{idVendor}=="0471", MODE="0666"" >> 99-android.rules
echo "#PMC-Sierra" >> 99-android.rules
echo "SUBSYSTEM=="usb", SYSFS{idVendor}=="04da", MODE="0666"" >> 99-android.rules
\end{verbatim}
\end{maxipage}
\begin{maxipage}
\begin{verbatim}
echo "#Qualcomm" >> 99-android.rules
echo "SUBSYSTEM=="usb", SYSFS{idVendor}=="05c6", MODE="0666"" >> 99-android.rules
echo "#SK Telesys" >> 99-android.rules
echo "SUBSYSTEM=="usb", SYSFS{idVendor}=="1f53", MODE="0666"" >> 99-android.rules
echo "#Samsung" >> 99-android.rules
echo "SUBSYSTEM=="usb", SYSFS{idVendor}=="04e8", MODE="0666"" >> 99-android.rules
echo "#Sharp" >> 99-android.rules
echo "SUBSYSTEM=="usb", SYSFS{idVendor}=="04dd", MODE="0666"" >> 99-android.rules
echo "#Sony Ericsson" >> 99-android.rules
echo "SUBSYSTEM=="usb", SYSFS{idVendor}=="0fce", MODE="0666"" >> 99-android.rules
echo "#Toshiba" >> 99-android.rules
echo "SUBSYSTEM=="usb", SYSFS{idVendor}=="0930", MODE="0666"" >> 99-android.rules
echo "#ZTE" >> 99-android.rules
echo "SUBSYSTEM=="usb", SYSFS{idVendor}=="19D2", MODE="0666"" >> 99-android.rules
sudo mv -f 99-android.rules /etc/udev/rules.d/
sudo chmod a+r /etc/udev/rules.d/99-android.rules
\end{verbatim}
\end{maxipage}

\section{Optional: Eclipse IDE}
One of the most widespread IDEs is eclipse. It is not limited to Android, and can be used for all targets. It is not required for \xc, but its installation is described here as an example. Eclipse is quite heavyweight, and many developers prefer other IDEs for \xc development.

To install, download the eclipse installer (Sometimes called "\emph{Ooomph!}" for some reason) from here:\\
\texttt{https://www.eclipse.org/downloads/}

Important: Install the CDT version of eclipse for C development, not the Android/Java package, even if you plan developing for Android. In addition, it is very convenient to install the git support (egit).

The current stable version is \emph{eclipse mars} (4.5) and works with OpenJDK 7 or 8, the new \emph{eclipse neon} 4.6, currently RC2, is also quite stable, and requires OpenJDK 8. Both can be installed with the installer.

You can also install the ADT (Android development tools) package for better integration with Android.

Next, create a new project, by generating a make project from existing sources files. Choose your xcsoar source directory which contains the makefile.

Important: After you have added the sources, eclipse will start indexing all files. If you have already started \texttt{make} before this time, then a lot of files have been downloaded for the various libraries which are exctracted/built within the \xc directory (most notably the boost libraries). Indexing all these takes a very long time, and a lot of heap space, so you should probably stop the indexer right away. In addition you should probably exclude these directories from the indexer for the future.

For this, in the C/C++ scope, right-click on the "output" directory in the file tree on the left side, select "Properties", then "Resource/Resource Filters" and add a filter. In the "add filter" dialog, choose "exclude all", "files and folders", "all children (recursive)" and set the Filter details to "Name matches *".
This will exclude the output tree from the indexer, leading to a minimal index.

\section{Build setup for the Manual}
To set up the LaTeX -- specifically, LuaLaTex -- system for compiling the manual, you need to install the following packages in addition to the steps from section \ref{sec:developsetup-base}, \ref{sec:developsetup-general} and \ref{sec:developsetup-sourcecode}:

\begin{verbatim}
sudo apt-get install texlive \
  texlive-latex-extra \
  texlive-luatex \
  texlive-lang-french \
  texlive-lang-polish \
  texlive-lang-german \
  liblocale-po-perl
\end{verbatim}

Now you can compile the manuals with 

\begin{verbatim}
make manual
\end{verbatim}

This creates a lot of warnings, but this is normal.

\section{Optional: modern LaTeX editor for editing the Manual}
Most people today edit LaTeX files in specific editors, as this is much more comfortable and efficient. This is highly recommended especially if you are not very familiar with LaTeX: learning it is very easy with a modern editor. Here, we install TeXstudio as an example, as it is very widespread and supports the rather rare LuaLaTeX well.

To install, get the relevant package:

\begin{verbatim}
sudo apt-get install texstudio
\end{verbatim}

As the directory tree of \xc is very unusual for a LaTeX project, we need to make some special configurations in order to allow for quick compiling from within the editor, and for full synctex functionality:

In "Options / Configure TeXStudio", enable "show advanced options".

In "Options / Configure TeXStudio / Commands / Commands / LuaLaTeX", replace
\begin{verbatim}
lualatex -synctex=1 -interaction=nonstopmode %.tex
\end{verbatim}
with
\begin{verbatim}
lualatex -synctex=1 -interaction=nonstopmode
 -output-directory=?a)../../../output/manual %.tex
\end{verbatim}

In "Options / Configure TexStudio / Build / Build Options / Addition Search Paths":\\
Enter in \emph{both} fields ("Log file" and in the field "PDF File"):
\begin{verbatim}
../../../output/manual/
\end{verbatim}


Add the following line to \emph{both} the \texttt{.profile} and the \texttt{.bashrc} file of your user directory:
\begin{maxipage}
\begin{verbatim}
export TEXINPUTS="..:../../../output/manual:../../../output/manual/en:../../..:"
\end{verbatim}
\end{maxipage}


Finally, you need to run "make manual" in the \xc base directory at least once from the command line before you can compile from within the TexStudio interface. This creates the path structure and generates the figure files which are included into the manual. Of course, if you change figures, you might have to run "make manual" again.

Inside TeXStudio, open the file "XCSoar-manual.tex" (or one of the other root files) and right-click on this file to "set as explicit root document", in the structure view on the left. Now you are good to go. Make changes and press F5 to see the result immediately.







\chapter{GNU General Public License}\label{cha:gnu-general-public}


Ceci est une traduction non officielle de la licence publique générale GNU (GNU GPL) en français (\url{http://www.fsffrance.org/gpl/gpl-fr.fr.html}). Elle n'a pas été publiée par la Free Software Foundation et n'établit pas juridiquement les termes de distribution des logiciels qui utilisent la GNU GPL – seul le texte anglais original de la GNU GPL le fait. Cependant, nous espérons que cette traduction aidera les francophones à mieux comprendre la GNU GPL.\\ 
Le texte original en anglais apparaît ensuite.

This is an unofficial translation of the GNU General Public License into French (\url{http://www.fsffrance.org/gpl/gpl-fr.fr.html}). It was not published by the Free Software Foundation, and does not legally state the distribution terms for software that uses the GNU GPL—only the original English text of the GNU GPL does that. However, we hope that this translation will help French speakers understand the GNU GPL better.\\
The original English text is included after.
\newpage

{\small
	
	\begin{center}
	
Licence Publique Générale GNU

	\end{center}

Les licences de la plupart des logiciels sont conçues pour vous enlever toute liberté de les partager et de les modifier.

A contrario, la Licence Publique Générale est destinée à garantir votre liberté de partager et de modifier les logiciels libres, et à assurer que ces logiciels soient libres pour tous leurs utilisateurs.

La présente Licence Publique Générale s'applique à la plupart des logiciels de la Free Software Foundation, ainsi qu'à tout autre programme pour lequel ses auteurs s'engagent à l'utiliser.

(Certains autres logiciels de la Free Software Foundation sont couverts par la GNU Lesser General Public License à la place.)

Vous pouvez aussi l'appliquer aux programmes qui sont les vôtres.

Quand nous parlons de logiciels libres, nous parlons de liberté, non de prix.

Nos licences publiques générales sont conçues pour vous donner l'assurance d'être libres de distribuer des copies des logiciels libres (et de facturer ce service, si vous le souhaitez), de recevoir le code source ou de pouvoir l'obtenir si vous le souhaitez, de pouvoir modifier les logiciels ou en utiliser des éléments dans de nouveaux programmes libres et de savoir que vous pouvez le faire.

Pour protéger vos droits, il nous est nécessaire d'imposer des limitations qui interdisent à quiconque de vous refuser ces droits ou de vous demander d'y renoncer.

Certaines responsabilités vous incombent en raison de ces limitations si vous distribuez des copies de ces logiciels, ou si vous les modifiez.

Par exemple, si vous distribuez des copies d'un tel programme, à titre gratuit ou contre une rémunération, vous devez accorder aux destinataires tous les droits dont vous disposez.

Vous devez vous assurer qu'eux aussi reçoivent ou puissent disposer du code source.

Et vous devez leur montrer les présentes conditions afin qu'ils aient connaissance de leurs droits.

Nous protégeons vos droits en deux étapes : (1) nous sommes titulaires des droits d'auteur du logiciel, et (2) nous vous délivrons cette licence, qui vous donne l'autorisation légale de copier, distribuer et/ou modifier le logiciel.

En outre, pour la protection de chaque auteur ainsi que la nôtre, nous voulons nous assurer que chacun comprenne que ce logiciel libre ne fait l'objet d'aucune garantie.

Si le logiciel est modifié par quelqu'un d'autre puis transmis à des tiers, nous voulons que les destinataires soient mis au courant que ce qu'ils ont reçu n'est pas le logiciel d'origine, de sorte que tout problème introduit par d'autres ne puisse entacher la réputation de l'auteur originel.

En définitive, un programme libre restera à la merci des brevets de logiciels.

Nous souhaitons éviter le risque que les redistributeurs d'un programme libre fassent des demandes individuelles de licence de brevet, ceci ayant pour effet de rendre le programme propriétaire.

Pour éviter cela, nous établissons clairement que toute licence de brevet doit être concédée de façon à ce que l'usage en soit libre pour tous ou bien qu'aucune licence ne soit concédée.

Les termes exacts et les conditions de copie, distribution et modification sont les suivants:
Conditions de copie, distribution et modification de la Licence Publique Générale GNU.

	\begin{enumerate}

\item La présente Licence s'applique à tout programme ou tout autre ouvrage contenant un avis, apposé par le titulaire des droits d'auteur, stipulant qu'il peut être distribué au titre des conditions de la présente Licence Publique Générale.

Ci-après, le "Programme" désigne l'un quelconque de ces programmes ou ouvrages, et un "ouvrage fondé sur le Programme" désigne soit le Programme, soit un ouvrage qui en dérive au titre des lois sur le droit d'auteur : en d'autres termes, un ouvrage contenant le Programme ou une partie de ce dernier, soit à l'identique, soit avec des modifications et/ou traduit dans un autre langage.

(Ci-après, le terme "modification" implique, sans s'y réduire, le terme traduction)

Chaque concessionaire sera désigné par "vous".

Les activités autres que la copie, la distribution et la modification ne sont pas couvertes par la présente Licence ; elles sont hors de son champ d'application.

L'opération consistant à exécuter le Programme n'est soumise à aucune limitation et les sorties du programme ne sont couvertes que si leur contenu constitue un ouvrage fondé sur le Programme (indépendamment du fait qu'il ait été réalisé par l'exécution du Programme).

La validité de ce qui précède dépend de ce que fait le Programme.

\item  Vous pouvez copier et distribuer des copies à l'identique du code source du Programme tel que vous l'avez reçu, sur n'importe quel support, du moment que vous apposiez sur chaque copie, de manière ad hoc et parfaitement visible, l'avis de droit d'auteur adéquat et une exonération de garantie ; que vous gardiez intacts tous les avis faisant référence à la présente Licence et à l'absence de toute garantie ; et que vous fournissiez à tout destinataire du Programme autre que vous-même un exemplaire de la présente Licence en même temps que le Programme.

Vous pouvez faire payer l'acte physique de transmission d'une copie, et vous pouvez, à votre discrétion, proposer une garantie contre rémunération.

\item  Vous pouvez modifier votre copie ou des copies du Programme ou n'importe quelle partie de celui-ci, créant ainsi un ouvrage fondé sur le Programme, et copier et distribuer de telles modifications ou ouvrage selon les termes de l'Article 1 ci-dessus, à condition de vous conformer également à chacune des obligations suivantes :


\begin{enumerate}
\item Vous devez munir les fichiers modifiés d'avis bien visibles stipulants que vous avez modifié ces fichiers, ainsi que la date de chaque modification ;

\item Vous devez prendre les dispositions nécessaires pour que tout ouvrage que vous distribuez ou publiez, et qui, en totalité ou en partie, contient ou est fondé sur le Programme - ou une partie quelconque de ce dernier - soit concédé comme un tout, à titre gratuit, à n'importe quel tiers, au titre des conditions de la présente Licence.

\item Si le programme modifié lit habituellement des instructions de façon interactive lorsqu'on l'exécute, vous devez, quand il commence son exécution pour ladite utilisation interactive de la manière la plus usuelle, faire en sorte qu'il imprime ou affiche une annonce comprenant un avis de droit d'auteur ad hoc, et un avis stipulant qu'il n'y a pas de garantie (ou bien indiquant que c'est vous qui fournissez la garantie), et que les utilisateurs peuvent redistribuer le programme en respectant les présentes obligations, et expliquant à l'utilisateur comment voir une copie de la présente Licence.


\end{enumerate}
(Exception : si le Programme est lui-même interactif mais n'imprime pas habituellement une telle annonce, votre ouvrage fondé sur le Programme n'est pas obligé d'imprimer une annonce).

Ces obligations s'appliquent à l'ouvrage modifié pris comme un tout.

Si des éléments identifiables de cet ouvrage ne sont pas fondés sur le Programme et peuvent raisonnablement être considérés comme des ouvrages indépendants distincts en eux mêmes, alors la présente Licence et ses conditions ne s'appliquent pas à ces éléments lorsque vous les distribuez en tant qu'ouvrages distincts.

Mais lorsque vous distribuez ces mêmes éléments comme partie d'un tout, lequel constitue un ouvrage fondé sur le Programme, la distribution de ce tout doit être soumise aux conditions de la présente Licence, et les autorisations qu'elle octroie aux autres concessionnaires s'étendent à l'ensemble de l'ouvrage et par conséquent à chaque et toute partie indifférement de qui l'a écrite.

Par conséquent, l'objet du présent article n'est pas de revendiquer des droits ou de contester vos droits sur un ouvrage entièrement écrit par vous; son objet est plutôt d'exercer le droit de contrôler la distribution d'ouvrages dérivés ou d'ouvrages collectifs fondés sur le Programme.

De plus, la simple proximité du Programme avec un autre ouvrage qui n'est pas fondé sur le Programme (ou un ouvrage fondé sur le Programme) sur une partition d'un espace de stockage ou un support de distribution ne place pas cet autre ouvrage dans le champ d'application de la présente Licence.

\item  Vous pouvez copier et distribuer le Programme (ou un ouvrage fondé sur lui, selon l'Article 2) sous forme de code objet ou d'exécutable, selon les termes des Articles 1 et 2 ci-dessus, à condition que vous accomplissiez l'un des points suivants :


\begin{enumerate}
\item L'accompagner de l'intégralité du code source correspondant, sous une forme lisible par un ordinateur, lequel doit être distribué au titre des termes des Articles 1 et 2 ci-dessus, sur un support habituellement utilisé pour l'échange de logiciels; ou,

\item L'accompagner d'une proposition écrite, valable pendant au moins trois ans, de fournir à tout tiers, à un tarif qui ne soit pas supérieur à ce que vous coûte l'acte physique de réaliser une distribution source, une copie intégrale du code source correspondant sous une forme lisible par un ordinateur, qui sera distribuée au titre des termes des Articles 1 et 2 ci-dessus, sur un support habituellement utilisé pour l'échange de logiciels; ou,

\item  L'accompagner des informations reçues par vous concernant la proposition de distribution du code source correspondant. (Cette solution n'est autorisée que dans le cas d'une distribution non commerciale et seulement si vous avez reçu le programme sous forme de code objet ou d'exécutable accompagné d'une telle proposition - en conformité avec le sous-Article b ci-dessus.)
\end{enumerate}

Le code source d'un ouvrage désigne la forme favorite pour travailler à des modifications de cet ouvrage. Pour un ouvrage exécutable, le code source intégral désigne la totalité du code source de la totalité des modules qu'il contient, ainsi que les éventuels fichiers de définition des interfaces qui y sont associés, ainsi que les scripts utilisés pour contrôler la compilation et l'installation de l'exécutable. Cependant, par exception spéciale, le code source distribué n'est pas censé inclure quoi que ce soit de normalement distribué (que ce soit sous forme source ou binaire) avec les composants principaux (compilateur, noyau, et autre) du système d'exploitation sur lequel l'exécutable tourne, à moins que ce composant lui-même n'accompagne l'exécutable.

Si distribuer un exécutable ou un code objet consiste à offrir un accès permettant leur copie depuis un endroit particulier, alors l'offre d'un accès équivalent pour copier le code source depuis le même endroit compte comme une distribution du code source - même si les tiers ne sont pas contraints de copier le source en même temps que le code objet.

\item  Vous ne pouvez copier, modifier, concéder en sous-licence, ou distribuer le Programme, sauf tel qu'expressément prévu par la présente Licence. Toute tentative de copier, modifier, concéder en sous-licence, ou distribuer le Programme d'une autre manière est réputée non valable, et met immédiatement fin à vos droits au titre de la présente Licence. Toutefois, les tiers ayant reçu de vous des copies, ou des droits, au titre de la présente Licence ne verront pas leurs autorisations résiliées aussi longtemps que ledits tiers se conforment pleinement à elle.

\item  Vous n'êtes pas obligé d'accepter la présente Licence étant donné que vous ne l'avez pas signée. Cependant, rien d'autre ne vous accorde l'autorisation de modifier ou distribuer le Programme ou les ouvrages fondés sur lui. Ces actions sont interdites par la loi si vous n'acceptez pas la présente Licence. En conséquence, en modifiant ou distribuant le Programme (ou un ouvrage quelconque fondé sur le Programme), vous signifiez votre acceptation de la présente Licence en le faisant, et de toutes ses conditions concernant la copie, la distribution ou la modification du Programme ou d'ouvrages fondés sur lui.

\item  Chaque fois que vous redistribuez le Programme (ou n'importe quel ouvrage fondé sur le Programme), une licence est automatiquement concédée au destinataire par le concédant originel de la licence, l'autorisant à copier, distribuer ou modifier le Programme, sous réserve des présentes conditions. Vous ne pouvez imposer une quelconque limitation supplémentaire à l'exercice des droits octroyés au titre des présentes par le destinataire. Vous n'avez pas la responsabilité d'imposer le respect de la présente Licence à des tiers.

\item  Si, conséquement à une décision de justice ou l'allégation d'une transgression de brevet ou pour toute autre raison (non limitée à un probleme de brevet), des obligations vous sont imposées (que ce soit par jugement, conciliation ou autre) qui contredisent les conditions de la présente Licence, elles ne vous excusent pas des conditions de la présente Licence. Si vous ne pouvez distribuer de manière à satisfaire simultanément vos obligations au titre de la présente Licence et toute autre obligation pertinente, alors il en découle que vous ne pouvez pas du tout distribuer le Programme. Par exemple, si une licence de brevet ne permettait pas une redistribution sans redevance du Programme par tous ceux qui reçoivent une copie directement ou indirectement par votre intermédiaire, alors la seule façon pour vous de satisfaire à la fois à la licence du brevet et à la présente Licence serait de vous abstenir totalement de toute distribution du Programme.

Si une partie quelconque de cet article est tenue pour nulle ou inopposable dans une circonstance particulière quelconque, l'intention est que le reste de l'article s'applique. La totalité de la section s'appliquera dans toutes les autres circonstances.

Cet article n'a pas pour but de vous induire à transgresser un quelconque brevet ou d'autres revendications à un droit de propriété ou à contester la validité de la moindre de ces revendications ; cet article a pour seul objectif de protéger l'intégrité du système de distribution du logiciel libre, qui est mis en oeuvre par la pratique des licenses publiques. De nombreuses personnes ont fait de généreuses contributions au large spectre de logiciels distribués par ce système en se fiant à l'application cohérente de ce système ; il appartient à chaque auteur/donateur de décider si il ou elle veut distribuer du logiciel par l'intermédiaire d'un quelconque autre système et un concessionaire ne peut imposer ce choix.

Cet article a pour but de rendre totalement limpide ce que l'on pense être une conséquence du reste de la présente Licence.

\item  Si la distribution et/ou l'utilisation du Programme est limitée dans certains pays que ce soit par des brevets ou par des interfaces soumises au droit d'auteur, le titulaire originel des droits d'auteur qui décide de couvrir le Programme par la présente Licence peut ajouter une limitation géographique de distribution explicite qui exclue ces pays afin que la distribution soit permise seulement dans ou entre les pays qui ne sont pas ainsi exclus. Dans ce cas, la présente Licence incorpore la limitation comme si elle était écrite dans le corps de la présente Licence.

\item  La Free Software Foundation peut, de temps à autre, publier des versions révisées et/ou nouvelles de la Licence Publique Générale. De telles nouvelles versions seront similaires à la présente version dans l'esprit mais pourront différer dans le détail pour prendre en compte de nouvelles problématiques ou inquiétudes.

Chaque version possède un numéro de version la distinguant. Si le Programme précise le numéro de version de la présente Licence qui s'y applique et "une version ultérieure quelconque", vous avez le choix de suivre les conditions de la présente version ou de toute autre version ultérieure publiée par la Free Software Foundation. Si le Programme ne spécifie aucun numéro de version de la présente Licence, vous pouvez choisir une version quelconque publiée par la Free Software Foundation à quelque moment que ce soit.

\item  Si vous souhaitez incorporer des parties du Programme dans d'autres programmes libres dont les conditions de distribution sont différentes, écrivez à l'auteur pour lui en demander l'autorisation. Pour les logiciels dont la Free Software Foundation est titulaire des droits d'auteur, écrivez à la Free Software Foundation ; nous faisons parfois des exceptions dans ce sens. Notre décision sera guidée par le double objectif de préserver le statut libre de tous les dérivés de nos logiciels libres et de promouvoir le partage et la réutilisation des logiciels en général.
ABSENCE DE GARANTIE

\item  COMME LA LICENCE DU PROGRAMME EST CONCEDEE A TITRE GRATUIT, AUCUNE GARANTIE NE S'APPLIQUE AU PROGRAMME, DANS LES LIMITES AUTORISEES PAR LA LOI APPLICABLE. SAUF MENTION CONTRAIRE ECRITE, LES TITULAIRES DU DROIT D'AUTEUR ET/OU LES AUTRES PARTIES FOURNISSENT LE PROGRAMME "EN L'ETAT", SANS AUCUNE GARANTIE DE QUELQUE NATURE QUE CE SOIT, EXPRESSE OU IMPLICITE, Y COMPRIS, MAIS SANS Y ETRE LIMITE, LES GARANTIES IMPLICITES DE COMMERCIABILITE ET DE LA CONFORMITE A UNE UTILISATION PARTICULIERE. VOUS ASSUMEZ LA TOTALITE DES RISQUES LIES A LA QUALITE ET AUX PERFORMANCES DU PROGRAMME. SI LE PROGRAMME SE REVELAIT DEFECTUEUX, LE COUT DE L'ENTRETIEN, DES REPARATIONS OU DES CORRECTIONS NECESSAIRES VOUS INCOMBENT INTEGRALEMENT.

\item  EN AUCUN CAS, SAUF LORSQUE LA LOI APPLICABLE OU UNE CONVENTION ECRITE L'EXIGE, UN TITULAIRE DE DROIT D'AUTEUR QUEL QU'IL SOIT, OU TOUTE PARTIE QUI POURRAIT MODIFIER ET/OU REDISTRIBUER LE PROGRAMME COMME PERMIS CI-DESSUS, NE POURRAIT ETRE TENU POUR RESPONSABLE A VOTRE EGARD DES DOMMAGES, INCLUANT LES DOMMAGES GENERIQUES, SPECIFIQUES, SECONDAIRES OU CONSECUTIFS, RESULTANT DE L'UTILISATION OU DE L'INCAPACITE D'UTILISER LE PROGRAMME (Y COMPRIS, MAIS SANS Y ETRE LIMITE, LA PERTE DE DONNEES, OU LE FAIT QUE DES DONNEES SOIENT RENDUES IMPRECISES, OU LES PERTES EPROUVEES PAR VOUS OU PAR DES TIERS, OU LE FAIT QUE LE PROGRAMME ECHOUE A INTEROPERER AVEC UN AUTRE PROGRAMME QUEL QU'IL SOIT) MEME SI LE DIT TITULAIRE DU DROIT D'AUTEUR OU LE PARTIE CONCERNEE A ETE AVERTI DE L'EVENTUALITE DE TELS DOMMAGES.

		\end{enumerate}


FIN DES CONDITIONS

\vspace{1cm}
Sur les pages suivantes, la version original en anglais de la General Public License. 
\newpage 

	\begin{center}
		GNU GENERAL PUBLIC LICENSE
		
		Version 2, June 1991
		
		
	\end{center}
	
	Copyright (C) 1989, 1991 Free Software Foundation, Inc.
	
	59 Temple Place, Suite 330, Boston, MA  02111-1307  USA
	
	Everyone is permitted to copy and distribute verbatim copies
	of this license document, but changing it is not allowed.
	
	\begin{center}
		Preamble
	\end{center}
	
	The licenses for most software are designed to take away your
	freedom to share and change it.  By contrast, the GNU General Public
	License is intended to guarantee your freedom to share and change free
	software--to make sure the software is free for all its users.  This
	General Public License applies to most of the Free Software
	Foundation's software and to any other program whose authors commit to
	using it.  (Some other Free Software Foundation software is covered by
	the GNU Library General Public License instead.)  You can apply it to
	your programs, too.
	
	When we speak of free software, we are referring to freedom, not
	price.  Our General Public Licenses are designed to make sure that you
	have the freedom to distribute copies of free software (and charge for
	this service if you wish), that you receive source code or can get it
	if you want it, that you can change the software or use pieces of it
	in new free programs; and that you know you can do these things.
	
	To protect your rights, we need to make restrictions that forbid
	anyone to deny you these rights or to ask you to surrender the rights.
	These restrictions translate to certain responsibilities for you if you
	distribute copies of the software, or if you modify it.
	
	For example, if you distribute copies of such a program, whether
	gratis or for a fee, you must give the recipients all the rights that
	you have.  You must make sure that they, too, receive or can get the
	source code.  And you must show them these terms so they know their
	rights.
	
	We protect your rights with two steps: (1) copyright the software, and
	(2) offer you this license which gives you legal permission to copy,
	distribute and/or modify the software.
	
	Also, for each author's protection and ours, we want to make certain
	that everyone understands that there is no warranty for this free
	software.  If the software is modified by someone else and passed on, we
	want its recipients to know that what they have is not the original, so
	that any problems introduced by others will not reflect on the original
	authors' reputations.
	
	Finally, any free program is threatened constantly by software
	patents.  We wish to avoid the danger that redistributors of a free
	program will individually obtain patent licenses, in effect making the
	program proprietary.  To prevent this, we have made it clear that any
	patent must be licensed for everyone's free use or not licensed at all.
	
	The precise terms and conditions for copying, distribution and
	modification follow.
	
	\begin{center}
		GNU GENERAL PUBLIC LICENSE
		
		TERMS AND CONDITIONS FOR COPYING, DISTRIBUTION AND MODIFICATION
	\end{center}
	
	\begin{enumerate}
		\item This License applies to any program or other work which contains
		a notice placed by the copyright holder saying it may be distributed
		under the terms of this General Public License.  The ``Program'', below,
		refers to any such program or work, and a ``work based on the Program''
		means either the Program or any derivative work under copyright law:
		that is to say, a work containing the Program or a portion of it,
		either verbatim or with modifications and/or translated into another
		language.  (Hereinafter, translation is included without limitation in
		the term ``modification''.)  Each licensee is addressed as ``you''.
		
		Activities other than copying, distribution and modification are not
		covered by this License; they are outside its scope.  The act of
		running the Program is not restricted, and the output from the Program
		is covered only if its contents constitute a work based on the
		Program (independent of having been made by running the Program).
		Whether that is true depends on what the Program does.
		
		\item You may copy and distribute verbatim copies of the Program's
		source code as you receive it, in any medium, provided that you
		conspicuously and appropriately publish on each copy an appropriate
		copyright notice and disclaimer of warranty; keep intact all the
		notices that refer to this License and to the absence of any warranty;
		and give any other recipients of the Program a copy of this License
		along with the Program.
		
		You may charge a fee for the physical act of transferring a copy, and
		you may at your option offer warranty protection in exchange for a fee.
		
		\item You may modify your copy or copies of the Program or any portion
		of it, thus forming a work based on the Program, and copy and
		distribute such modifications or work under the terms of Section 1
		above, provided that you also meet all of these conditions:
		
		\begin{enumerate}
			\item You must cause the modified files to carry prominent notices
			stating that you changed the files and the date of any change.
			
			\item You must cause any work that you distribute or publish, that in
			whole or in part contains or is derived from the Program or any
			part thereof, to be licensed as a whole at no charge to all third
			parties under the terms of this License.
			
			\item If the modified program normally reads commands interactively
			when run, you must cause it, when started running for such
			interactive use in the most ordinary way, to print or display an
			announcement including an appropriate copyright notice and a
			notice that there is no warranty (or else, saying that you provide
			a warranty) and that users may redistribute the program under
			these conditions, and telling the user how to view a copy of this
			License.  (Exception: if the Program itself is interactive but
			does not normally print such an announcement, your work based on
			the Program is not required to print an announcement.)
		\end{enumerate}
		
		These requirements apply to the modified work as a whole.  If
		identifiable sections of that work are not derived from the Program,
		and can be reasonably considered independent and separate works in
		themselves, then this License, and its terms, do not apply to those
		sections when you distribute them as separate works.  But when you
		distribute the same sections as part of a whole which is a work based
		on the Program, the distribution of the whole must be on the terms of
		this License, whose permissions for other licensees extend to the
		entire whole, and thus to each and every part regardless of who wrote it.
		
		Thus, it is not the intent of this section to claim rights or contest
		your rights to work written entirely by you; rather, the intent is to
		exercise the right to control the distribution of derivative or
		collective works based on the Program.
		
		In addition, mere aggregation of another work not based on the Program
		with the Program (or with a work based on the Program) on a volume of
		a storage or distribution medium does not bring the other work under
		the scope of this License.
		
		\item You may copy and distribute the Program (or a work based on it,
		under Section 2) in object code or executable form under the terms of
		Sections 1 and 2 above provided that you also do one of the following:
		
		\begin{enumerate}
			\item Accompany it with the complete corresponding machine-readable
			source code, which must be distributed under the terms of Sections
			1 and 2 above on a medium customarily used for software interchange; or,
			
			\item Accompany it with a written offer, valid for at least three
			years, to give any third party, for a charge no more than your
			cost of physically performing source distribution, a complete
			machine-readable copy of the corresponding source code, to be
			distributed under the terms of Sections 1 and 2 above on a medium
			customarily used for software interchange; or,
			
			\item Accompany it with the information you received as to the offer
			to distribute corresponding source code.  (This alternative is
			allowed only for noncommercial distribution and only if you
			received the program in object code or executable form with such
			an offer, in accord with Subsection b above.)
		\end{enumerate}
		
		The source code for a work means the preferred form of the work for
		making modifications to it.  For an executable work, complete source
		code means all the source code for all modules it contains, plus any
		associated interface definition files, plus the scripts used to
		control compilation and installation of the executable.  However, as a
		special exception, the source code distributed need not include
		anything that is normally distributed (in either source or binary
		form) with the major components (compiler, kernel, and so on) of the
		operating system on which the executable runs, unless that component
		itself accompanies the executable.
		
		If distribution of executable or object code is made by offering
		access to copy from a designated place, then offering equivalent
		access to copy the source code from the same place counts as
		distribution of the source code, even though third parties are not
		compelled to copy the source along with the object code.
		
		\item You may not copy, modify, sublicense, or distribute the Program
		except as expressly provided under this License.  Any attempt
		otherwise to copy, modify, sublicense or distribute the Program is
		void, and will automatically terminate your rights under this License.
		However, parties who have received copies, or rights, from you under
		this License will not have their licenses terminated so long as such
		parties remain in full compliance.
		
		\item You are not required to accept this License, since you have not
		signed it.  However, nothing else grants you permission to modify or
		distribute the Program or its derivative works.  These actions are
		prohibited by law if you do not accept this License.  Therefore, by
		modifying or distributing the Program (or any work based on the
		Program), you indicate your acceptance of this License to do so, and
		all its terms and conditions for copying, distributing or modifying
		the Program or works based on it.
		
		\item Each time you redistribute the Program (or any work based on the
		Program), the recipient automatically receives a license from the
		original licensor to copy, distribute or modify the Program subject to
		these terms and conditions.  You may not impose any further
		restrictions on the recipients' exercise of the rights granted herein.
		You are not responsible for enforcing compliance by third parties to
		this License.
		
		\item If, as a consequence of a court judgment or allegation of patent
		infringement or for any other reason (not limited to patent issues),
		conditions are imposed on you (whether by court order, agreement or
		otherwise) that contradict the conditions of this License, they do not
		excuse you from the conditions of this License.  If you cannot
		distribute so as to satisfy simultaneously your obligations under this
		License and any other pertinent obligations, then as a consequence you
		may not distribute the Program at all.  For example, if a patent
		license would not permit royalty-free redistribution of the Program by
		all those who receive copies directly or indirectly through you, then
		the only way you could satisfy both it and this License would be to
		refrain entirely from distribution of the Program.
		
		If any portion of this section is held invalid or unenforceable under
		any particular circumstance, the balance of the section is intended to
		apply and the section as a whole is intended to apply in other
		circumstances.
		
		It is not the purpose of this section to induce you to infringe any
		patents or other property right claims or to contest validity of any
		such claims; this section has the sole purpose of protecting the
		integrity of the free software distribution system, which is
		implemented by public license practices.  Many people have made
		generous contributions to the wide range of software distributed
		through that system in reliance on consistent application of that
		system; it is up to the author/donor to decide if he or she is willing
		to distribute software through any other system and a licensee cannot
		impose that choice.
		
		This section is intended to make thoroughly clear what is believed to
		be a consequence of the rest of this License.
		
		\item If the distribution and/or use of the Program is restricted in
		certain countries either by patents or by copyrighted interfaces, the
		original copyright holder who places the Program under this License
		may add an explicit geographical distribution limitation excluding
		those countries, so that distribution is permitted only in or among
		countries not thus excluded.  In such case, this License incorporates
		the limitation as if written in the body of this License.
		
		\item The Free Software Foundation may publish revised and/or new versions
		of the General Public License from time to time.  Such new versions will
		be similar in spirit to the present version, but may differ in detail to
		address new problems or concerns.
		
		Each version is given a distinguishing version number.  If the Program
		specifies a version number of this License which applies to it and ``any
		later version'', you have the option of following the terms and conditions
		either of that version or of any later version published by the Free
		Software Foundation.  If the Program does not specify a version number of
		this License, you may choose any version ever published by the Free Software
		Foundation.
		
		\item If you wish to incorporate parts of the Program into other free
		programs whose distribution conditions are different, write to the author
		to ask for permission.  For software which is copyrighted by the Free
		Software Foundation, write to the Free Software Foundation; we sometimes
		make exceptions for this.  Our decision will be guided by the two goals
		of preserving the free status of all derivatives of our free software and
		of promoting the sharing and reuse of software generally.
		
	\end{enumerate}
	
	\subsection*{No warranty}
	
	BECAUSE THE PROGRAM IS LICENSED FREE OF CHARGE, THERE IS NO WARRANTY
	FOR THE PROGRAM, TO THE EXTENT PERMITTED BY APPLICABLE LAW.  EXCEPT
	WHEN OTHERWISE STATED IN WRITING THE COPYRIGHT HOLDERS AND/OR OTHER
	PARTIES PROVIDE THE PROGRAM ``AS IS'' WITHOUT WARRANTY OF ANY KIND,
	EITHER EXPRESSED OR IMPLIED, INCLUDING, BUT NOT LIMITED TO, THE
	IMPLIED WARRANTIES OF MERCHANTABILITY AND FITNESS FOR A PARTICULAR
	PURPOSE.  THE ENTIRE RISK AS TO THE QUALITY AND PERFORMANCE OF THE
	PROGRAM IS WITH YOU.  SHOULD THE PROGRAM PROVE DEFECTIVE, YOU ASSUME
	THE COST OF ALL NECESSARY SERVICING, REPAIR OR CORRECTION.
	
	IN NO EVENT UNLESS REQUIRED BY APPLICABLE LAW OR AGREED TO IN WRITING
	WILL ANY COPYRIGHT HOLDER, OR ANY OTHER PARTY WHO MAY MODIFY AND/OR
	REDISTRIBUTE THE PROGRAM AS PERMITTED ABOVE, BE LIABLE TO YOU FOR
	DAMAGES, INCLUDING ANY GENERAL, SPECIAL, INCIDENTAL OR CONSEQUENTIAL
	DAMAGES ARISING OUT OF THE USE OR INABILITY TO USE THE PROGRAM
	(INCLUDING BUT NOT LIMITED TO LOSS OF DATA OR DATA BEING RENDERED
	INACCURATE OR LOSSES SUSTAINED BY YOU OR THIRD PARTIES OR A FAILURE
	OF THE PROGRAM TO OPERATE WITH ANY OTHER PROGRAMS), EVEN IF SUCH
	HOLDER OR OTHER PARTY HAS BEEN ADVISED OF THE POSSIBILITY OF SUCH
	DAMAGES.
	
}
\endinput


\end{document}
